\input{preamble}

\begin{document}

\title{Homological algebra}
\maketitle

\phantomsection
\label{section-phantom}

\tableofcontents

\section{Categories of G-modules and
R-modules}
\label{section-g-and-r-modules}

\noindent
Fix a group $G$ and a ring $R$. We write
$\mathit{Mod}_R$ for the category of left
$R$-modules and $R$-linear maps.

\begin{definition}
\label{definition-g-module}
A {\it $G$-module} is an abelian group $A$
with a left action of $G$ by group
automorphisms. Morphisms of $G$-modules are
$G$-equivariant group homomorphisms. Their
category is denoted $\mathit{Mod}_G$.
\end{definition}

\noindent
Equivalently, a $G$-module is a left
$\mathbb{Z}[G]$-module. The convention in
Definition \ref{definition-g-module} agrees
with the Stacks Project,
\href{https://stacks.math.columbia.edu/tag/%
04JP}{Tag 04JP}.

\begin{definition}
\label{definition-invariants-coinvariants}
Let $A$ be a $G$-module. The {\it
invariants} and {\it coinvariants} of $A$
are, respectively,
$$
A^G
=
\{a\in A\mid ga=a\quad\forall g\in G\}
$$
and
$$
A_G
=
A\Big/\sum_{g\in G}(g-1)A.
$$
\end{definition}

\noindent
The invariants in Definition
\ref{definition-invariants-coinvariants}
are the elements already fixed by $G$.
Taking coinvariants forces $ga$ and $a$
to become equal. See the
Stacks Project,
\href{https://stacks.math.columbia.edu/tag/%
0A2H}{Tag 0A2H} and
\href{https://stacks.math.columbia.edu/tag/%
09WT}{Tag 09WT}.

\begin{example}
\label{example-invariants-quotient}
Let $G=\{1,s\}$ be the group of order two.
Let $A=\mathbb{Z}$, with $sa=-a$, and let
$B=2\mathbb{Z}$. Then
$$
A^G=B^G=0,
$$
while $G$ acts trivially on
$A/B\cong\mathbb{Z}/2\mathbb{Z}$. Hence
$$
(A/B)^G\ne A^G/B^G.
$$
\end{example}

\noindent
Example \ref{example-invariants-quotient}
makes the following comment from the
lecture concrete.

\begin{quote}
``Understanding natural operations
(lie $(-)^G$), applied to natural
constructions
(like quotient)
and the failure of the most naive
guess (namely $(A/B)^G=A^G/B^G$)
is kind of what homological algebra
is about.''
\end{quote}

\section{Preadditive and additive categories}
\label{section-additive-categories}

\noindent
{\bf Note.} Apart from the standard
definition, I highlight from Jethro's
lecture that he proved 
Lemma \ref{lemma-preadditive-direct-sum}.


\medskip\noindent
Here is the definition of a preadditive
category.

\begin{definition}
\label{definition-preadditive}
A category $\mathcal{A}$ is called {\it
preadditive} if each
morphism set $\Mor_\mathcal{A}(x, y)$ is
endowed
with the structure of an abelian group such
that the
compositions
$$
\Mor(x, y) \times \Mor(y, z)
\longrightarrow
\Mor(x, z)
$$
are bilinear. A functor $F : \mathcal{A} \to
\mathcal{B}$ of
preadditive categories is called {\it
additive} if and only
if $F : \Mor(x, y) \to \Mor(F(x), F(y))$
is a homomorphism of abelian groups for all
$x, y \in \Ob(\mathcal{A})$.
\end{definition}

\noindent
In particular for every $x, y$ there exists
at least
one morphism $x \to y$, namely the zero map.

\begin{lemma}
\label{lemma-preadditive-zero}
Let $\mathcal{A}$ be a preadditive category.
Let $x$ be an object of $\mathcal{A}$.
The following are equivalent
\begin{enumerate}
\item $x$ is an initial object,
\item $x$ is a final object, and
\item $\text{id}_x = 0$ in
$\Mor_\mathcal{A}(x, x)$.
\end{enumerate}
Furthermore, if such an object $0$ exists,
then a morphism
$\alpha : y \to z$ factors through $0$ if and
only if $\alpha = 0$.
\end{lemma}

\begin{proof}
First assume that $x$ is either (1) initial
or (2) final.
In both cases, it follows that $\Mor(x,x)$ is
a trivial abelian group
containing $\text{id}_x$, thus $\text{id}_x =
0$ in
$\Mor(x, x)$, which shows that each of (1)
and (2) implies (3).

\medskip\noindent
Now assume that $\text{id}_x = 0$ in
$\Mor(x,x)$. Let $y$
be an arbitrary object of $\mathcal{A}$ and
let $f \in \Mor(x ,y)$.
Denote $C : \Mor(x,x) \times \Mor(x,y) \to
\Mor(x,y)$ the composition map.
Then $f = C(0, f)$ and since $C$ is bilinear
we have $C(0, f) = 0$.
Thus $f = 0$. Hence $x$ is initial in
$\mathcal{A}$.
A similar argument for $f \in \Mor(y, x)$ can
be used to show that
$x$ is also final. Thus (3) implies both (1)
and (2).
\end{proof}

\begin{definition}
\label{definition-zero-object}
In a preadditive category $\mathcal{A}$, we
call an object that is both
final and initial (as in Lemma
\ref{lemma-preadditive-zero}) a
{\it zero object} and denote it $0$.
\end{definition}

\begin{lemma}
\label{lemma-preadditive-direct-sum}
Let $\mathcal{A}$ be a preadditive category.
Let $x, y \in \Ob(\mathcal{A})$.
If the product $x \times y$ exists, then so
does
the coproduct $x \amalg y$.
If the coproduct $x \amalg y$ exists, then so
does
the product $x \times y$. In this case
also $x \amalg y \cong x \times y$.
\end{lemma}

\begin{proof}
Suppose that $z = x \times y$ with
projections
$p : z \to x$ and $q : z \to y$. Denote $i :
x \to z$
the morphism corresponding to $(1, 0)$.
Denote $j : y \to z$
the morphism corresponding to $(0, 1)$. Thus
we have the
commutative diagram
$$
\xymatrix{
x \ar[rr]^1 \ar[rd]^i & & x \\
& z \ar[ru]^p \ar[rd]^q & \\
y \ar[rr]^1 \ar[ru]^j & & y
}
$$
where the diagonal compositions are zero. It
follows that
$i \circ p + j \circ q : z \to z$ is the
identity since
it is a morphism which upon composing with
$p$ gives $p$
and upon composing with $q$ gives $q$.
Suppose given morphisms $a : x \to w$ and $b
: y \to w$.
Then we can form the map $a \circ p + b \circ
q : z \to w$.
In this way we get a bijection $\Mor(z, w)
= \Mor(x, w) \times \Mor(y, w)$ which
show that $z = x \amalg y$.

\medskip\noindent
We leave it to the reader to construct the
morphisms
$p, q$ given a coproduct $x \amalg y$ instead
of a
product.
\end{proof}

\begin{definition}
\label{definition-direct-sum}
Given a pair of objects $x, y$ in a
preadditive category $\mathcal{A}$,
the {\it direct sum} $x \oplus y$ of $x$ and
$y$ is the direct
product $x \times y$ endowed with the
morphisms
$i, j, p, q$ as in Lemma
\ref{lemma-preadditive-direct-sum} above.
\end{definition}

\begin{remark}
\label{remark-direct-sum}
Note that the proof of Lemma
\ref{lemma-preadditive-direct-sum}
shows that given $p$ and $q$ the morphisms
$i$, $j$ are uniquely
determined by the rules $p \circ i =
\text{id}_x$,
$q \circ j = \text{id}_y$, $p \circ j = 0$,
$q \circ i = 0$.
Moreover, we automatically have
$i \circ p + j \circ q = \text{id}_{x \oplus
y}$.
Similarly, given $i$, $j$ the morphisms $p$
and $q$ are uniquely determined.
Finally, given objects $x, y, z$ and
morphisms
$i : x \to z$, $j : y \to z$, $p : z \to x$
and
$q : z \to y$ such that $p \circ i =
\text{id}_x$,
$q \circ j = \text{id}_y$, $p \circ j = 0$,
$q \circ i = 0$
and $i \circ p + j \circ q = \text{id}_z$,
then $z$
is the direct sum of $x$ and $y$ with the
four morphisms
equal to $i, j, p, q$.
\end{remark}

\begin{lemma}
\label{lemma-additive-additive}
Let $\mathcal{A}$, $\mathcal{B}$ be
preadditive categories.
Let $F : \mathcal{A} \to \mathcal{B}$ be an
additive functor.
Then $F$ transforms direct sums to direct
sums and zero to zero.
\end{lemma}

\begin{proof}
Suppose $F$ is additive. A direct sum $z$
of $x$ and $y$ is characterized by having
morphisms
$i : x \to z$, $j : y \to z$, $p : z \to x$
and
$q : z \to y$ such that $p \circ i =
\text{id}_x$,
$q \circ j = \text{id}_y$, $p \circ j = 0$,
$q \circ i = 0$
and $i \circ p + j \circ q = \text{id}_z$,
according
to Remark \ref{remark-direct-sum}. Clearly
$F(x), F(y), F(z)$
and the morphisms $F(i), F(j), F(p), F(q)$
satisfy exactly the
same relations (by additivity) and we see
that $F(z)$ is
a direct sum of $F(x)$ and $F(y)$.
Hence, $F$ transforms direct sums to direct
sums.

\medskip\noindent
To see that $F$ transforms zero to zero, use
the
characterization (3) of the zero object in
Lemma \ref{lemma-preadditive-zero}.
\end{proof}

\begin{definition}
\label{definition-additive-category}
A category $\mathcal{A}$ is called {\it
additive}
if it is preadditive and finite products
exist, in other
words it has a zero object and direct sums.
\end{definition}

\noindent
Namely the empty product is a finite product
and
if it exists, then it is a final object.

\begin{definition}
\label{definition-kernel}
Let $\mathcal{A}$ be a preadditive category.
Let $f : x \to y$ be a morphism.
\begin{enumerate}
\item A {\it kernel} of $f$ is a morphism
$i : z \to x$ such that (a) $f \circ i = 0$
and (b)
for any $i' : z' \to x$ such that $f \circ i'
= 0$ there
exists a unique morphism $g : z' \to z$ such
that
$i' = i \circ g$.
\item If the kernel of $f$ exists, then we
denote
this $\Ker(f) \to x$.
\item A {\it cokernel} of $f$ is a morphism
$p : y \to z$ such that (a) $p \circ f = 0$
and (b)
for any $p' : y \to z'$ such that $p' \circ f
= 0$ there
exists a unique morphism $g : z \to z'$ such
that
$p' = g \circ p$.
\item If a cokernel of $f$ exists we denote
this
$y \to \Coker(f)$.
\item If a kernel of $f$ exists, then a {\it
coimage
of $f$} is a cokernel for the morphism
$\Ker(f) \to x$.
\item If a kernel and coimage exist then we
denote this
$x \to \Coim(f)$.
\item If a cokernel of $f$ exists, then the
{\it image of
$f$} is a kernel of the morphism $y \to
\Coker(f)$.
\item If a cokernel and image of $f$ exist
then we denote
this $\Im(f) \to y$.
\end{enumerate}
\end{definition}

\noindent
In the above definition, we have spoken of
``the kernel'' and
``the cokernel'', tacitly using their
uniqueness
up to unique isomorphism. This follows from
the Yoneda lemma
(Categories, Section
\ref{categories-section-opposite})
because the kernel of $f : x \to y$
represents the functor
sending an object $z$ to the set
$\Ker(\Mor_\mathcal{A}(z, x) \to
\Mor_\mathcal{A}(z, y))$.
The case of cokernels is dual.

\medskip\noindent
We first relate the direct sum to kernels as
follows.

\begin{lemma}
\label{lemma-additive-cat-biproduct-kernel}
Let $\mathcal{C}$ be a preadditive category.
Let $x \oplus y$ with morphisms $i, j, p, q$
as in
Lemma \ref{lemma-preadditive-direct-sum}
be a direct sum in $\mathcal{C}$. Then $i : x
\to x \oplus y$
is a kernel of $q : x \oplus y \rightarrow
y$. Dually, $p$ is
a cokernel for $j$.
\end{lemma}

\begin{proof}
Let $f : z' \to x \oplus y$ be a morphism
such that $q \circ f = 0$.
We have to show that there exists a unique
morphism $g : z' \to x$
such that $f = i \circ g$. Since $i \circ p +
j \circ q$ is the identity on
$x \oplus y$ we see that
$$
f = (i \circ p + j \circ q) \circ f = i \circ
p \circ f
$$
and hence $g = p \circ f$ works. Uniqueness
holds because $p \circ i$
is the identity on $x$. The proof of the
second statement is dual.
\end{proof}

\begin{lemma}
\label{lemma-kernel-mono}
Let $\mathcal{C}$ be a preadditive category.
Let $f : x \to y$ be a morphism in
$\mathcal{C}$.
\begin{enumerate}
\item If a kernel of $f$ exists, then
this kernel is a monomorphism.
\item If a cokernel of $f$ exists, then
this cokernel is an epimorphism.
\item If a kernel and coimage of $f$ exist,
then
the coimage is an epimorphism.
\item If a cokernel and image of $f$ exist,
then
the image is a monomorphism.
\end{enumerate}
\end{lemma}

\begin{proof}
Part (1) follows easily from the uniqueness
required in the
definition of a kernel. The proof of (2) is
dual.
Part (3) follows from (2), since the coimage
is a cokernel.
Similarly, (4) follows from (1).
\end{proof}

\begin{lemma}
\label{lemma-coim-im-map}
Let $f : x \to y$ be a morphism in a
preadditive category
such that the kernel, cokernel, image and
coimage all exist.
Then $f$ can be factored uniquely as
$x \to \Coim(f) \to \Im(f) \to y$.
\end{lemma}

\begin{proof}
There is a canonical morphism $\Coim(f) \to
y$
because $\Ker(f) \to x \to y$ is zero.
The composition $\Coim(f) \to y \to
\Coker(f)$
is zero, because it is the unique morphism
which gives
rise to the morphism $x \to y \to \Coker(f)$
which
is zero
(the uniqueness follows from
Lemma \ref{lemma-kernel-mono} (3)).
Hence $\Coim(f) \to y$ factors uniquely
through
$\Im(f) \to y$, which gives us the desired
map.
\end{proof}

\begin{example}
\label{example-not-abelian}
Let $k$ be a field.
Consider the category
of filtered vector spaces over $k$.
(See Definition \ref{definition-filtered}.)
Consider the filtered vector spaces $(V, F)$
and $(W, F)$ with
$V = W = k$ and
$$
F^iV
=
\left\{
\begin{matrix}
V & \text{if} & i < 0 \\
0 & \text{if} & i \geq 0
\end{matrix}
\right.
\text{ and }
F^iW
=
\left\{
\begin{matrix}
W & \text{if} & i \leq 0 \\
0 & \text{if} & i > 0
\end{matrix}
\right.
$$
The map $f : V \to W$ corresponding to
$\text{id}_k$ on the underlying
vector spaces has trivial kernel and cokernel
but is not
an isomorphism. Note also that $\Coim(f) = V$
and $\Im(f) = W$.
This means that the category of filtered
vector spaces over $k$
is not abelian.
\end{example}





\section{Abelian categories}
\label{section-abelian-categories}

\noindent
An abelian category is a category satisfying
just enough axioms so the
snake lemma holds. An axiom (that is
sometimes forgotten)
is that the canonical map $\Coim(f) \to
\Im(f)$
of Lemma \ref{lemma-coim-im-map} is always an
isomorphism.
Example \ref{example-not-abelian} shows that
it is necessary.

Also see Remark \ref{remark-snake-lemma-les}
on how the Snake Lemma is used
to prove the existence of the
induced long exact sequence
in homology.

\begin{definition}
\label{definition-abelian-category}
A category $\mathcal{A}$ is {\it abelian} if
it is additive, if all kernels and cokernels
exist,
and if the natural map $\Coim(f) \to \Im(f)$
is an isomorphism for all morphisms $f$ of
$\mathcal{A}$.
\end{definition}

\begin{lemma}
\label{lemma-abelian-opposite}
Let $\mathcal{A}$ be a preadditive category.
The additions on sets of morphisms make
$\mathcal{A}^{opp}$ into a preadditive
category.
Furthermore, $\mathcal{A}$ is additive if and
only if $\mathcal{A}^{opp}$
is additive, and
$\mathcal{A}$ is abelian if and only if
$\mathcal{A}^{opp}$ is abelian.
\end{lemma}

\begin{proof}
The first statement is straightforward.
To see that $\mathcal{A}$ is additive if and
only if $\mathcal{A}^{opp}$
is additive, recall that additivity can be
characterized by
the existence of a zero object and direct
sums, which are both
preserved when passing to the opposite
category.
Finally, to see that
$\mathcal{A}$ is abelian if and only if
$\mathcal{A}^{opp}$ is abelian,
observes that kernels, cokernels, images and
coimages in
$\mathcal{A}^{opp}$ correspond to
cokernels, kernels, coimages and images in
$\mathcal{A}$,
respectively.
\end{proof}

\begin{definition}
\label{definition-injective-surjective}
Let $f : x \to y$ be a morphism in an abelian
category.
\begin{enumerate}
\item We say $f$ is {\it injective} if
$\Ker(f) = 0$.
\item We say $f$ is {\it surjective} if
$\Coker(f) = 0$.
\end{enumerate}
If $x \to y$ is injective, then we say that
$x$ is a {\it subobject}
of $y$ and we use the notation $x \subset y$.
If $x \to y$ is
surjective, then we say that $y$ is a {\it
quotient} of $x$.
\end{definition}

\begin{lemma}
\label{lemma-characterize-injective}
Let $f : x \to y$ be a morphism in an abelian
category $\mathcal{A}$. Then
\begin{enumerate}
\item $f$ is injective if and only if $f$ is
a monomorphism,
\item $f$ is surjective if and only if $f$ is
an epimorphism, and
\item $f$ is an isomorphism if and only if
$f$ is injective and surjective.
\end{enumerate}
\end{lemma}

\begin{proof}
Proof of (1). Recall that $\Ker(f)$ is an
object representing the
functor sending $z$ to
$\Ker(\Mor_\mathcal{A}(z, x) \to
\Mor_\mathcal{A}(z, y))$, see
Definition \ref{definition-kernel}.
Thus $\Ker(f)$ is $0$ if and only if
$\Mor_\mathcal{A}(z, x) \to
\Mor_\mathcal{A}(z, y)$
is injective for all $z$ if and only if $f$
is a monomorphism.
The proof of (2) is similar.
For the proof of (3) note that an isomorphism
is both a monomorphism
and epimorphism, which by (1), (2) proves $f$
is injective
and surjective. If $f$ is both injective and
surjective, then
$x = \Coim(f)$ and $y = \Im(f)$ whence $f$ is
an isomorphism.
\end{proof}

\noindent
In an abelian category, if $x \subset y$ is a
subobject,
then we denote
$$
y/x = \Coker(x \to y).
$$

\begin{lemma}
\label{lemma-colimit-abelian-category}
Let $\mathcal{A}$ be an abelian category.
All finite limits and finite colimits exist
in $\mathcal{A}$.
\end{lemma}

\begin{proof}
To show that finite limits exist it suffices
to show
that finite products and equalizers exist,
see
Categories, Lemma
\ref{categories-lemma-finite-limits-exist}.
Finite products exist
by definition and the equalizer of $a, b : x
\to y$ is
the kernel of $a - b$. The argument for
finite colimits
is similar but dual to this.
\end{proof}

\begin{example}
\label{example-fibre-product-pushouts}
Let $\mathcal{A}$ be an abelian category.
Pushouts and fibre products in $\mathcal{A}$
have the following
simple descriptions:
\begin{enumerate}
\item If $a : x \to y$, $b : z \to y$ are
morphisms in $\mathcal{A}$, then
we have the fibre product:
$x \times_y z = \Ker((a, -b) : x \oplus z \to
y)$.
\item If $a : y \to x$, $b : y \to z$ are
morphisms in $\mathcal{A}$, then
we have the pushout:
$x \amalg_y z = \Coker((a, -b) : y \to x
\oplus z)$.
\end{enumerate}
\end{example}

\begin{definition}
\label{definition-exact}
Let $\mathcal{A}$ be an additive category.
Consider a sequence of morphisms
$$
\ldots \to x \to y \to z \to \ldots
\quad\text{or}\quad
x_1 \to x_2 \to \ldots \to x_n
$$
in $\mathcal{A}$. We say such a sequence is a
{\it complex} if the
composition of any two consecutive (drawn)
arrows is zero.
If $\mathcal{A}$ is abelian then we say a
complex of the first
type above is {\it exact at $y$} if $\Im(x
\to y) = \Ker(y \to z)$
and we say a complex of the second kind is
{\it exact at $x_i$}
where $1 < i < n$ if
$\Im(x_{i - 1} \to x_i) = \Ker(x_i \to x_{i +
1})$. We a
sequence as above is {\it exact} or is an
{\it exact sequence} or is an
{\it exact complex} if it is a complex and
exact at every object (in
the first case) or exact at $x_i$ for all $1
< i < n$ (in the second case).
There are variants of these notions for
sequences of the form
$$
\ldots \to x_{-3} \to x_{-2} \to x_{-1}
\quad\text{and}\quad
x_1 \to x_2 \to x_3 \to \ldots
$$
A {\it short exact sequence} is an exact
complex of the form
$$
0 \to A  \to B \to C \to 0.
$$
\end{definition}

\noindent
In the following lemma we assume the reader
knows what it means
for a sequence of abelian groups to be exact.

\begin{lemma}
\label{lemma-check-exactness}
Let $\mathcal{A}$ be an abelian category.
Let $0 \to M_1 \to M_2 \to M_3 \to 0$ be a
complex of $\mathcal{A}$.
\begin{enumerate}
\item $M_1 \to M_2 \to M_3 \to 0$ is exact if
and only if
$$
0 \to \Hom_\mathcal{A}(M_3, N) \to
\Hom_\mathcal{A}(M_2, N) \to
\Hom_\mathcal{A}(M_1, N)
$$
is an exact sequence of abelian groups for
all objects $N$ of
$\mathcal{A}$, and
\item $0 \to M_1 \to M_2 \to M_3$ is exact if
and only if
$$
0 \to \Hom_\mathcal{A}(N, M_1) \to
\Hom_\mathcal{A}(N, M_2) \to
\Hom_\mathcal{A}(N, M_3)
$$
is an exact sequence of abelian groups for
all objects $N$ of $\mathcal{A}$.
\end{enumerate}
\end{lemma}

\begin{proof}
Omitted. Hint: See
Algebra, Lemma \ref{algebra-lemma-hom-exact}.
\end{proof}

\begin{definition}
\label{definition-ses-split}
Let $\mathcal{A}$ be an abelian category.
Let $i : A \to B$ and $q : B \to C$ be
morphisms
of $\mathcal{A}$ such that
$0 \to A \to B \to C \to 0$ is a short
exact sequence. We say the short exact
sequence is {\it split} if there exist
morphisms $j : C \to B$ and $p : B \to A$
such
that $(B, i, j, p, q)$ is the direct sum of
$A$ and $C$.
\end{definition}

\begin{lemma}
\label{lemma-ses-split}
Let $\mathcal{A}$ be an abelian category.
Let $0 \to A \to B \to C \to 0$
be a short exact sequence.
\begin{enumerate}
\item Given a morphism $s : C \to B$ right
inverse to
$B \to C$, there exists a unique $\pi : B \to
A$
such that $(s, \pi)$ splits the short exact
sequence
as in Definition \ref{definition-ses-split}.
\item Given a morphism $\pi : B \to A$ left
inverse to
$A \to B$, there exists a unique $s : C \to
B$
such that $(s, \pi)$ splits the short exact
sequence
as in Definition \ref{definition-ses-split}.
\end{enumerate}
\end{lemma}

\begin{proof}
Omitted.
\end{proof}

\begin{lemma}
\label{lemma-characterize-cartesian}
Let $\mathcal{A}$ be an abelian category. Let
$$
\xymatrix{
w\ar[r]^f\ar[d]_g
& y\ar[d]^h\\
x\ar[r]^k
& z
}
$$
be a commutative diagram.
\begin{enumerate}
\item The diagram is cartesian if and only if
$$
0 \to w \xrightarrow{(g, f)} x \oplus y
\xrightarrow{(k, -h)} z
$$
is exact.
\item The diagram is cocartesian if and only
if
$$
w \xrightarrow{(g, -f)} x \oplus y
\xrightarrow{(k, h)} z \to 0
$$
is exact.
\end{enumerate}
\end{lemma}

\begin{proof}
Let $u = (g, f) : w \to x \oplus y$ and $v =
(k, -h) : x \oplus y \to z$.
Let $p : x \oplus y \to x$ and $q : x \oplus
y \to y$ be the canonical
projections. Let $i : \Ker(v) \to x \oplus y$
be the canonical
injection. By Example
\ref{example-fibre-product-pushouts}, the
diagram is
cartesian if and only if there exists an
isomorphism
$r : \Ker(v) \to w$ with $f \circ r = q \circ
i$ and
$g \circ r = p \circ i$. The sequence
$0 \to w \overset{u} \to x \oplus y
\overset{v} \to z$ is exact if and
only if there exists an isomorphism $r :
\Ker(v) \to w$ with
$u \circ r = i$. But given $r : \Ker(v) \to
w$, we have
$f \circ r = q \circ i$ and $g \circ r = p
\circ i$ if and
only if $q \circ u \circ r= f \circ r = q
\circ i$ and
$p \circ u \circ r = g \circ r = p \circ i$,
hence if and only if
$u \circ r = i$. This proves (1), and then
(2) follows by duality.
\end{proof}

\begin{lemma}
\label{lemma-cartesian-kernel}
Let $\mathcal{A}$ be an abelian category. Let
$$
\xymatrix{
w\ar[r]^f\ar[d]_g
& y\ar[d]^h\\
x\ar[r]^k
& z
}
$$
be a commutative diagram.
\begin{enumerate}
\item If the diagram is cartesian, then the
morphism
$\Ker(f)\to\Ker(k)$ induced by $g$ is an
isomorphism.
\item If the diagram is cocartesian, then the
morphism
$\Coker(f)\to\Coker(k)$ induced by $h$ is an
isomorphism.
\end{enumerate}
\end{lemma}

\begin{proof}
Suppose the diagram is cartesian. Let
$e:\Ker(f)\to\Ker(k)$ be induced by $g$. Let
$i:\Ker(f)\to w$ and $j:\Ker(k)\to x$ be the
canonical
injections. There exists $t:\Ker(k)\to w$
with $f\circ t=0$
and $g\circ t=j$. Hence, there exists
$u:\Ker(k)\to\Ker(f)$
with $i\circ u=t$. It follows
$g\circ i\circ u\circ e=g\circ t\circ
e=j\circ e=g\circ i$ and
$f\circ i\circ u\circ e=0=f\circ i$, hence
$i\circ u\circ e=i$. Since
$i$ is a monomorphism this implies $u\circ
e=\text{id}_{\Ker(f)}$.
Furthermore, we have $j\circ e\circ u=g\circ
i\circ u=g\circ t=j$.
Since $j$ is a monomorphism this implies
$e\circ u=\text{id}_{\Ker(k)}$.
This proves (1). Now, (2) follows by duality.
\end{proof}

\begin{lemma}
\label{lemma-cartesian-cocartesian}
Let $\mathcal{A}$ be an abelian category. Let
$$
\xymatrix{
w\ar[r]^f\ar[d]_g
& y\ar[d]^h\\
x\ar[r]^k
& z
}
$$
be a commutative diagram.
\begin{enumerate}
\item If the diagram is cartesian and $k$ is
an epimorphism,
then the diagram is cocartesian and $f$ is an
epimorphism.
\item If the diagram is cocartesian and $g$
is a monomorphism,
then the diagram is cartesian and $h$ is a
monomorphism.
\end{enumerate}
\end{lemma}

\begin{proof}
Suppose the diagram is cartesian and $k$ is
an epimorphism.
Let $u = (g, f) : w \to x \oplus y$ and let
$v = (k, -h) : x \oplus y \to z$.
As $k$ is an epimorphism, $v$ is an
epimorphism, too. Therefore
and by Lemma
\ref{lemma-characterize-cartesian}, the
sequence
$0\to w\overset{u}\to x\oplus y\overset{v}\to
z\to 0$ is exact. Thus, the
diagram is cocartesian by Lemma
\ref{lemma-characterize-cartesian}. Finally,
$f$ is an epimorphism by Lemma
\ref{lemma-cartesian-kernel} and
Lemma \ref{lemma-characterize-injective}.
This proves (1), and (2)
follows by duality.
\end{proof}

\begin{lemma}
\label{lemma-epimo-unive-abeli-categ}
Let $\mathcal{A}$ be an abelian category.
\begin{enumerate}
\item If $x \to y$ is surjective, then for
every $z \to y$ the
projection $x \times_y z \to z$ is
surjective.
\item If $x \to y$ is injective, then for
every $x \to z$ the
morphism $z \to z \amalg_x y$ is injective.
\end{enumerate}
\end{lemma}

\begin{proof}
Immediately from Lemma
\ref{lemma-characterize-injective} and
Lemma \ref{lemma-cartesian-cocartesian}.
\end{proof}

\begin{lemma}
\label{lemma-check-exactness-fibre-product}
Let $\mathcal{A}$ be an abelian category. Let
$f:x\to y$ and $g:y\to z$
be morphisms with $g\circ f=0$. Then, the
following statements are equivalent:
\begin{enumerate}
\item The sequence $x\overset{f}\to
y\overset{g}\to z$ is exact.
\item For every $h:w\to y$ with $g\circ h=0$
there exist an object $v$,
an epimorphism $k:v\to w$ and a morphism
$l:v\to x$ with $h\circ k=f\circ l$.
\end{enumerate}
\end{lemma}

\begin{proof}
Let $i:\Ker(g)\to y$ be the canonical
injection. Let
$p:x\to\Coim(f)$ be the canonical projection.
Let
$j:\Im(f)\to\Ker(g)$ be the canonical
injection.

\medskip\noindent
Suppose (1) holds. Let $h:w\to y$ with
$g\circ h=0$. There exists
$c:w\to\Ker(g)$ with $i\circ c=h$.
Let $v=x\times_{\Ker(g)}w$ with canonical
projections
$k:v\to w$ and $l:v\to x$, so that $c\circ
k=j\circ p\circ l$.
Then, $h\circ k=i\circ c\circ k=i\circ j\circ
p\circ l=f\circ l$.
As $j\circ p$ is an epimorphism by
hypothesis, $k$ is an
epimorphism by Lemma
\ref{lemma-cartesian-cocartesian}. This
implies (2).

\medskip\noindent
Suppose (2) holds. Then, $g\circ i=0$. So,
there are an object
$w$, an epimorphism $k:w\to\Ker(g)$ and a
morphism
$l:w\to x$ with $f\circ l=i\circ k$. It
follows
$i\circ j\circ p\circ l=f\circ l=i\circ k$.
Since $i$ is a
monomorphism we see that $j\circ p\circ l=k$
is an epimorphism.
So, $j$ is an epimorphisms and thus an
isomorphism. This implies (1).
\end{proof}

\begin{lemma}
\label{lemma-exact-kernel-sequence}
Let $\mathcal{A}$ be an abelian category. Let
$$
\xymatrix{
x \ar[r]^f \ar[d]^\alpha &
y \ar[r]^g \ar[d]^\beta &
z \ar[d]^\gamma\\
u \ar[r]^k & v \ar[r]^l & w
}
$$
be a commutative diagram.
\begin{enumerate}
\item If the first row is exact and $k$ is a
monomorphism, then the induced
sequence $\Ker(\alpha) \to \Ker(\beta) \to
\Ker(\gamma)$
is exact.
\item If the second row is exact and $g$ is
an epimorphism, then the induced
sequence
$\Coker(\alpha) \to \Coker(\beta) \to
\Coker(\gamma)$
is exact.
\end{enumerate}
\end{lemma}

\begin{proof}
Suppose the first row is exact and $k$ is a
monomorphism. Let
$a:\Ker(\alpha)\to\Ker(\beta)$ and
$b:\Ker(\beta)\to\Ker(\gamma)$ be the induced
morphisms.
Let $h:\Ker(\alpha)\to x$, $i:\Ker(\beta)\to
y$ and
$j:\Ker(\gamma)\to z$ be the canonical
injections. As $j$ is
a monomorphism we have $b\circ a=0$. Let
$c:s\to\Ker(\beta)$
with $b\circ c=0$. Then, $g\circ i\circ
c=j\circ b\circ c=0$. By
Lemma
\ref{lemma-check-exactness-fibre-product}
there are an object $t$, an
epimorphism $d:t\to s$ and a morphism $e:t\to
x$ with
$i\circ c\circ d=f\circ e$. Then,
$k\circ \alpha\circ e=\beta\circ f\circ
e=\beta\circ i\circ c\circ d=0$.
As $k$ is a monomorphism we get $\alpha\circ
e=0$. So, there exists
$m:t\to\Ker(\alpha)$ with $h\circ m=e$. It
follows
$i\circ a\circ m=f\circ h\circ m=f\circ
e=i\circ c\circ d$.
As $i$ is a monomorphism we get $a\circ
m=c\circ d$. Thus,
Lemma
\ref{lemma-check-exactness-fibre-product}
implies (1), and then
(2) follows by duality.
\end{proof}

\noindent
(From Jethro's lecture.)
The point of the snake lemma is the
existence of the connecting morphism.
In particular, this allows us to
construct the long exact sequence in
cohomology from a short exact sequence;
that is, not only do we get induced
morphisms but also the weird connecting
morphism.

\begin{lemma}[Snake lemma, Jethro's version]
\label{lemma-snake-jethro}
In an abelian category $\mathcal{C}$,
let the following diagram be given,
with exact rows:
$$
\xymatrix{
& A \ar[r]^f \ar[d]^\alpha &
B \ar[r]^g \ar[d]^\beta &
C \ar[r] \ar[d]^\gamma &
0 \\
0 \ar[r] & A' \ar[r]^k & B' \ar[r]^l & C'
}
$$
then there is an exact sequence
$$
\xymatrix{
\ker(f)\ar[r]
&\ker(\alpha) \ar[r]
&\ker(\beta)\ar[r]
&\ker(\gamma)\ar[lld]\\
&\operatorname{coker}(\alpha)\ar[r]
&\operatorname{coker}(\beta)\ar[r]
&\operatorname{coker}(\gamma)\ar[r]
&\operatorname{coker}(g').
}
$$
\end{lemma}

\begin{proof}
Mostly routine, but exhausting.
For instance, one can check from
universal properties, why there
exists a morphism […]

The most interesting part is the 
morphism
$$
\ker(\gamma)\xrightarrow{\delta}
\operatorname{coker}(\alpha).
$$
Let's construct it. […]
\end{proof}

\begin{lemma}[Snake lemma, Stacks Project]
\label{lemma-snake}
Let $\mathcal{A}$ be an abelian category. Let
$$
\xymatrix{
& x \ar[r]^f \ar[d]^\alpha &
y \ar[r]^g \ar[d]^\beta &
z \ar[r] \ar[d]^\gamma &
0 \\
0 \ar[r] & u \ar[r]^k & v \ar[r]^l & w
}
$$
be a commutative diagram with exact rows.
\begin{enumerate}
\item There exists a unique morphism
$\delta : \Ker(\gamma) \to \Coker(\alpha)$
such that the diagram
$$
\xymatrix{
y \ar[d]_\beta &
y \times_z \Ker(\gamma) \ar[l]_{\pi'}
\ar[r]^{\pi} &
\Ker(\gamma) \ar[d]^\delta \\
v \ar[r]^{\iota'} & \Coker(\alpha) \amalg_u v
&
\Coker(\alpha) \ar[l]_\iota
}
$$
commutes, where $\pi$ and $\pi'$ are the
canonical projections
and $\iota$ and $\iota'$ are the canonical
coprojections.
\item The induced sequence
$$
\Ker(\alpha) \xrightarrow{f'} \Ker(\beta)
\xrightarrow{g'}
\Ker(\gamma) \xrightarrow{\delta}
\Coker(\alpha) \xrightarrow{k'}
\Coker(\beta) \xrightarrow{l'} \Coker(\gamma)
$$
is exact. If $f$ is injective then so is
$f'$, and if $l$ is
surjective then so is $l'$.
\end{enumerate}
\end{lemma}

\begin{proof}
As $\pi$ is an epimorphism and $\iota$ is a
monomorphism by
Lemma \ref{lemma-cartesian-cocartesian},
uniqueness of $\delta$ is clear.
Let $p = y \times_z \Ker(\gamma)$ and $q =
\Coker(\alpha) \amalg_u v$.
Let $h : \Ker(\beta) \to y$, $i :
\Ker(\gamma) \to z$ and
$j : \Ker(\pi) \to p$ be the canonical
injections.
Let $\pi'' : u \to \Coker(\alpha)$ be the
canonical projection.
Keeping in mind Lemma
\ref{lemma-cartesian-cocartesian} we get a
commutative
diagram with exact rows
$$
\xymatrix{
0 \ar[r] &
\Ker(\pi) \ar[r]^j &
p \ar[r]^{\pi} \ar[d]_{\pi'} &
\Ker(\gamma) \ar[d]_i \ar[r] & 0 \\
& x \ar[r]^f \ar[d]_\alpha & y \ar[r]^g
\ar[d]_\beta &
z \ar[d]_\gamma \ar[r] & 0 \\
0 \ar[r] & u \ar[r]^k \ar[d]_{\pi''} &
v \ar[r]^l \ar[d]_{\iota'} & w & \\
0 \ar[r] & \Coker(\alpha) \ar[r]^\iota & q &
&
}
$$
As $l \circ \beta \circ \pi' = \gamma \circ i
\circ \pi = 0$ and as the third
row of the diagram above is exact, there is
an $a : p \to u$
with $k \circ a = \beta \circ \pi'$. As the
upper right quadrangle of the
diagram above is cartesian, Lemma
\ref{lemma-cartesian-kernel} yields an
epimorphism $b : x \to \Ker(\pi)$ with $\pi'
\circ j \circ b = f$.
It follows
$k \circ a \circ j \circ b = \beta \circ \pi'
\circ j \circ b =
\beta \circ f = k \circ \alpha$.
As $k$ is a monomorphism this implies $a
\circ j \circ b = \alpha$. It follows
$\pi'' \circ a \circ j \circ b = \pi'' \circ
\alpha = 0$. As $b$ is an
epimorphism this
implies $\pi'' \circ a \circ j = 0$.
Therefore, as the top row of the diagram
above is exact, there exists
$\delta : \Ker(\gamma) \to \Coker(\alpha)$
with
$\delta \circ \pi = \pi'' \circ a$. It
follows
$\iota \circ \delta \circ \pi = \iota \circ
\pi'' \circ a =
\iota' \circ k \circ a = \iota' \circ \beta
\circ \pi'$
as desired.

\medskip\noindent
As the upper right quadrangle in the diagram
above is cartesian there
is a $c : \Ker(\beta) \to p$ with $\pi' \circ
c = h$ and $\pi \circ c = g'$.
It follows
$\iota \circ \delta \circ g' = \iota \circ
\delta \circ \pi \circ c =
\iota' \circ \beta \circ \pi' \circ c =
\iota' \circ \beta \circ h = 0$.
As $\iota$ is a monomorphism this implies
$\delta \circ g' = 0$.

\medskip\noindent
Next, let $d : r \to \Ker(\gamma)$ with
$\delta \circ d = 0$. Applying
Lemma
\ref{lemma-check-exactness-fibre-product}
to the exact sequence
$p \xrightarrow{\pi} \Ker(\gamma) \to 0$ and
$d$ yields an object $s$,
an epimorphism $m : s \to r$ and a morphism
$n : s \to p$ with
$\pi \circ n = d \circ m$. As
$\pi'' \circ a \circ n = \delta \circ d \circ
m = 0$,
applying Lemma
\ref{lemma-check-exactness-fibre-product} to
the exact sequence
$x \xrightarrow{\alpha} u \xrightarrow{p}
\Coker(\alpha)$ and
$a \circ n$ yields an object $t$, an
epimorphism $\varepsilon : t \to s$ and
a morphism $\zeta : t \to x$ with
$a \circ n \circ \varepsilon = \alpha \circ
\zeta$.
It holds
$\beta \circ \pi' \circ n \circ \varepsilon =
k \circ \alpha \circ \zeta = \beta \circ f
\circ \zeta$.
Let $\eta = \pi' \circ n \circ \varepsilon -
f \circ \zeta : t \to y$. Then,
$\beta \circ \eta = 0$. It follows that there
is a
$\vartheta : t \to \Ker(\beta)$ with $\eta =
h \circ \vartheta$. It holds
$i \circ g' \circ \vartheta = g \circ h \circ
\vartheta =
g \circ \pi' \circ n \circ \varepsilon - g
\circ f \circ \zeta =
i \circ \pi \circ n \circ \varepsilon = i
\circ d \circ m \circ \varepsilon$.
As $i$ is a monomorphism we get
$g' \circ \vartheta = d \circ m \circ
\varepsilon$.
Thus, as $m \circ \varepsilon$ is an
epimorphism,
Lemma
\ref{lemma-check-exactness-fibre-product}
implies that
$\Ker(\beta) \xrightarrow{g'} \Ker(\gamma)
\xrightarrow{\delta} \Coker(\alpha)$
is exact. Then, the claim follows by Lemma
\ref{lemma-exact-kernel-sequence}
and duality.
\end{proof}

\begin{lemma}
\label{lemma-snake-natural}
Let $\mathcal{A}$ be an abelian category. Let
$$
\xymatrix{
& & & x\ar[ld]\ar[rr]\ar[dd]^(.4)\alpha
& & y\ar[ld]\ar[rr]\ar[dd]^(.4)\beta
& & z\ar[ld]\ar[rr]\ar[dd]^(.4)\gamma
& & 0\\
& & x'\ar[rr]\ar[dd]^(.4){\alpha'}
& & y'\ar[rr]\ar[dd]^(.4){\beta'}
& & z'\ar[rr]\ar[dd]^(.4){\gamma'}
& & 0
& \\
& 0\ar[rr]
& & u\ar[ld]\ar[rr]
& & v\ar[ld]\ar[rr]
& & w\ar[ld]
& & \\
0\ar[rr]
& & u'\ar[rr]
& & v'\ar[rr]
& & w'
& & &
}
$$
be a commutative diagram with exact rows.
Then, the induced diagram
$$
\xymatrix@C=15pt{
\Ker(\alpha) \ar[r] \ar[d] &
\Ker(\beta) \ar[r] \ar[d] &
\Ker(\gamma) \ar[r]^(.45){\delta} \ar[d] &
\Coker(\alpha) \ar[r] \ar[d] &
\Coker(\beta) \ar[r] \ar[d] &
\Coker(\gamma) \ar[d] \\
\Ker(\alpha') \ar[r] &
\Ker(\beta') \ar[r] &
\Ker(\gamma') \ar[r]^(.45){\delta'} &
\Coker(\alpha') \ar[r] &
\Coker(\beta') \ar[r] &
\Coker(\gamma')
}
$$
commutes.
\end{lemma}

\begin{proof}
Omitted.
\end{proof}

\begin{lemma}
\label{lemma-four-lemma}
Let $\mathcal{A}$ be an abelian category. Let
$$
\xymatrix{
w \ar[r] \ar[d]^\alpha & x \ar[r]
\ar[d]^\beta & y \ar[r] \ar[d]^\gamma &
z \ar[d]^\delta \\
w' \ar[r] & x' \ar[r] & y' \ar[r] & z'
}
$$
be a commutative diagram with exact rows.
\begin{enumerate}
\item If $\alpha, \gamma$ are surjective and
$\delta$ is injective, then
$\beta$ is surjective.
\item If $\beta, \delta$ are injective and
$\alpha$ is surjective, then
$\gamma$ is injective.
\end{enumerate}
\end{lemma}

\begin{proof}
Assume $\alpha, \gamma$ are surjective and
$\delta$ is injective.
We may replace $w'$ by $\Im(w' \to x')$,
i.e., we may assume
that $w' \to x'$ is injective.
We may replace $z$ by $\Im(y \to z)$, i.e.,
we may assume that
$y \to z$ is surjective. Then we may apply
Lemma \ref{lemma-snake}
to
$$
\xymatrix{
& \Ker(y \to z) \ar[r] \ar[d] & y \ar[r]
\ar[d] & z \ar[r] \ar[d] & 0 \\
0 \ar[r] & \Ker(y' \to z') \ar[r] & y' \ar[r]
& z'
}
$$
to conclude that $\Ker(y \to z) \to \Ker(y'
\to z')$ is
surjective. Finally, we apply
Lemma \ref{lemma-snake}
to
$$
\xymatrix{
& w \ar[r] \ar[d] & x \ar[r] \ar[d] & \Ker(y
\to z) \ar[r] \ar[d] & 0 \\
0 \ar[r] & w' \ar[r] & x' \ar[r] & \Ker(y'
\to z')
}
$$
to conclude that $x \to x'$ is surjective.
This proves (1). The proof
of (2) is dual to this.
\end{proof}

\begin{lemma}
\label{lemma-five-lemma}
\begin{reference}
\cite[Lemma 4.5 page 16]{Eilenberg-Steenrod}
\end{reference}
Let $\mathcal{A}$ be an abelian category. Let
$$
\xymatrix{
v \ar[r] \ar[d]^\alpha &
w \ar[r] \ar[d]^\beta &
x \ar[r] \ar[d]^\gamma &
y \ar[r] \ar[d]^\delta &
z \ar[d]^\epsilon \\
v' \ar[r] & w' \ar[r] & x' \ar[r] & y' \ar[r]
& z'
}
$$
be a commutative diagram with exact rows. If
$\beta, \delta$
are isomorphisms, $\epsilon$ is injective,
and $\alpha$ is surjective
then $\gamma$ is an isomorphism.
\end{lemma}

\begin{proof}
Immediate consequence of
Lemma \ref{lemma-four-lemma}.
\end{proof}

\medskip\noindent
From Jethro's lecture:

\begin{example}[Non-abelian categories]
\label{example-non-abelian-categories}
\begin{enumerate}
\item Fix $M$ a smooth manifold,
and consider the category of finite
dimensional vector bundles.
This is not abelian.
Let $f$ be a map of the total spaces
of two vector bundles which is linear
on every fiber, i.e.
$$
\xymatrix{
E_1\ar[rr]^{f}\ar[dr]_{p_1}
&&E_2\ar[dl]^{p_2}\\
&M
}
$$
then ker and coker don't exist in general,
because the dimension of the fibers
might jump.

\item $M=\mathbb{R}$, $E=\mathbb{R}\times
V$,
$V=\mathbb{C}$ (so one-dimensional).

Let
$$
\xymatrix{
E \ar[dr]\ar[rr]&& E \ar[dl]\\
& \mathbb{R}
}
$$
with $f(v)=zv$, $z $ coordinate on
$\mathbb{R}$.
Then $\ker(f)=(\mathbb{R}\setminus 0)
\times \coprod 0 \times \mathbb{C}$.

The notion of coherent sheaf
is motivated by this example.

\item Other silly examples like the category
of all free abelian groups:
$\mathbb{Z} \xrightarrow{3}\mathbb{Z}$
has no cokernel.
\end{enumerate}
\end{example}

\medskip\noindent
We finish with the example of filtered
vector spaces mentioned in the lecture
(this was done by Jethro but the following
notes are AI-written).

\begin{definition}
\label{definition-filtered}
Let $k$ be a field. A {\it decreasing
filtration} on a vector space $V$ over $k$
is a collection of vector subspaces $F^pV$,
indexed by $p \in \mathbb{Z}$, such that
$$
F^{p + 1}V \subset F^pV
$$
for every $p$. A {\it filtered vector space}
is a pair $(V, F)$ consisting of a vector
space and a decreasing filtration on it.
A {\it morphism of filtered vector spaces}
from $(V, F)$ to $(W, G)$ is a linear map
$f : V \to W$ such that
$$
f(F^pV) \subset G^pW
$$
for every $p$. Such a morphism is called
{\it strict} if
$$
f(F^pV)=f(V)\cap G^pW
$$
for every $p$. We denote the resulting
category by $\mathit{Fil}(\mathit{Vec}_k)$.
No boundedness, exhaustiveness, or
separatedness condition is imposed on the
filtrations.
\end{definition}

\begin{lemma}
\label{lemma-filtered-vector-spaces}
The category $\mathit{Fil}(\mathit{Vec}_k)$
is additive and has all kernels and
cokernels. More precisely, let
$f : (V, F) \to (W, G)$ be a morphism, and
let $K = \Ker(f)$ and $Q = \Coker(f)$ be
the kernel and cokernel of the underlying
linear map. Their filtrations are
$$
F^pK=K\cap F^pV
$$
and
$$
G^pQ=(G^pW+f(V))/f(V).
$$
The canonical morphism
$$
\Coim(f)\longrightarrow\Im(f)
$$
is an isomorphism if and only if $f$ is
strict. Consequently,
$\mathit{Fil}(\mathit{Vec}_k)$ is not an
abelian category.
\end{lemma}

\begin{proof}
Addition of morphisms is inherited from
linear maps, and composition is bilinear.
The zero vector space with its unique
filtration is a zero object. Direct sums
are formed with the filtration
$$
H^p(V\oplus W)=F^pV\oplus G^pW.
$$
Hence the category is additive. The usual
kernel map is
filtration preserving for the filtration on
$K$ displayed above. If a morphism of
filtered vector spaces has zero composition
with $f$, its factorization through $K$ is
filtration preserving. Thus this is the
kernel in $\mathit{Fil}(\mathit{Vec}_k)$.

The filtration on $Q$ is the image of
$G^pW$ under the quotient map $W \to Q$.
The usual factorization through $Q$ of a
filtration preserving map that annihilates
$f(V)$ is again filtration preserving. Thus
this is the cokernel.

The filtration on the coimage corresponds,
under the usual vector space isomorphism
from the coimage to $f(V)$, to the subspaces
$f(F^pV)$. The filtration on the image is
given by the subspaces
$$
f(V)\cap G^pW.
$$
Hence the canonical morphism from the
coimage to the image is an isomorphism
exactly when $f$ is strict. The morphism in
Example \ref{example-not-abelian} is not
strict, so the category is not abelian.
\end{proof}

\section{Tensor product}
\label{section-tensor-product}

If the direct sum $A \oplus B$
'categorifies' $+$, surely what
categorifies $\cdot$ is $\otimes$,
the tensor product. Recall
that for vector spaces $U, V$
with bases $\{u_i\}, \{v_j\}$,
then $U \otimes V$ has basis
$\{u_i \otimes v_j\}$
(and has dimension $\dim(U \otimes V)
=\dim(U)\dim(V)$).

Universal property? The idea is that
a morphism $A \otimes B \to X$
is the same as a bilinear map
$A \times B \to X$.

$A \otimes B$ is caracterized by
the existence of a bijection
$\tilde{f} \leftrightarrow f$
of $\tilde{f}:A \otimes B \to X$,
and $f \in \operatorname{Bil}_{A,B}(X)$
such that… (exercise: such that what?).

\medskip\noindent
Another perspective is: for $T=A \otimes B$
fixed, we want the functos
$$
\Hom_{\underline{\operatorname{Ab}}}
(T,-):\underline{\operatorname{Ab}}
\to \underline{\operatorname{Set}}
$$
and
$$
\operatorname{Bil}_{A,B}(-):
\underline{\operatorname{Ab}}
\to \underline{\operatorname{Set}}
$$
to be isomorphic, i.e. that there exists
an natural isomorphism (i.e. invertible
natural transformation) between them.
That is:

\begin{definition}
\label{definition-tensor-product-categorical}
A tensor product of $A,B \in \underline{\operatorname{Ab}}$
is an object $T \in \underline{\operatorname{Ab}}$
and a natural isomorphism
$\theta:\Hom_{\operatorname{Ab}}(T,-)
\to \operatorname{Bil}_{A,B}(-)$
of functors $\underline{\operatorname{Ab}}
\to \underline{\operatorname{Set}}$.
\end{definition}

\noindent
If you decode the fancy categorical
language, it's the usual universal
property of $A \otimes B$.

\begin{lemma}
\label{lemma-tensor-product-uniqueness}
Two tensor products are naturally
isomorphic.
\end{lemma}

\noindent
Does a $\otimes$ exist in
$\underline{\operatorname{Ab}}$? Yes.
Let $A,B \in \underline{\operatorname{Ab}}$.
Consider $\mathbb{Z}[A \times B]$,
the free Abelian group on Cartesian
product. Let $J \subset \mathbb{Z}[A\times
B]$ be the submodule generated
by [distributivity of operations],
then $A \otimes B=\mathbb{Z}[A\times B]/J$
is a tensor product of $A$ and $B$.

\section{Complexes}
\label{section-complexes}

\noindent
The upshot here is:
complexes on an additive category
$\mathcal{A}$ form a category called
the {\it category of chain complexes of
$\mathcal{A}$},
denoted $\text{Ch}(\mathcal{A})$.
If $\mathcal{A}$ is abelian, then
$\operatorname{Ch}(\mathcal{A})$
is abelian too
(Lemma \ref{lemma-cat-chain-abelian},
by snake lemma).

For any $i \in \mathbf{Z}$ the $i$-th
homology group is a functor
$$
H_i : \text{Ch}(\mathcal{A}) \longrightarrow
\mathcal{A},
$$
which assigns to a chain complex $A_\bullet$
in an abelian category
its  $i$th {\it homology group}, defined by
the following formula
$$
H_i(A_\bullet) = \Ker(d_i)/\Im(d_{i + 1}).
$$
If $f : A_\bullet \to B_\bullet$ is a
morphism of chain
complexes of $\mathcal{A}$ then we get an
induced
morphism $H_i(f) : H_i(A_\bullet) \to
H_i(B_\bullet)$
because clearly
$f_i(\Ker(d_i : A_i \to A_{i - 1})) \subset
\Ker(d_i : B_i \to B_{i - 1})$, and similarly
for $\Im(d_{i + 1})$.
This proves that $H_i$ is indeed a functor.

Similarly,
for any $i \in \mathbf{Z}$ the $i$th {\it
cohomology group}
of a cochain complex $A^\bullet$ is defined
by
the following formula
$$
H^i(A^\bullet) = \Ker(d^i)/\Im(d^{i - 1}).
$$
If $f : A^\bullet \to B^\bullet$ is a
morphism of cochain
complexes of $\mathcal{A}$ then we get an
induced
morphism $H^i(f) : H^i(A^\bullet) \to
H^i(B^\bullet)$
because clearly
$f^i(\Ker(d^i : A^i \to A^{i + 1})) \subset
\Ker(d^i : B^i \to B^{i + 1})$, and similarly
for $\Im(d^{i - 1})$.
Thus we obtain a functor
$$
H^i : \text{CoCh}(\mathcal{A})
\longrightarrow \mathcal{A}.
$$




\medskip\noindent
Of course the notions of a chain complex and
a cochain complex
are dual and you only have to read one of the
two parts of
this section. So pick the one you like.
(Actually, this doesn't
quite work right since the conventions on
numbering things
are not adapted to an easy transition between
chain and cochain
complexes.)

\medskip\noindent
A {\it chain complex $A_\bullet$} in an
additive category $\mathcal{A}$
is a complex
$$
\ldots \to
A_{n + 1} \xrightarrow{d_{n + 1}}
A_n \xrightarrow{d_n}
A_{n - 1} \to
\ldots
$$
of $\mathcal{A}$. In other words, we are
given an object $A_i$ of
$\mathcal{A}$ for all $i \in \mathbf{Z}$ and
for
all $i \in \mathbf{Z}$ a morphism $d_i : A_i
\to A_{i - 1}$ such that
$d_{i - 1} \circ d_i = 0$ for all $i$. A {\it
morphism of chain
complexes $f : A_\bullet \to B_\bullet$} is
given by a
family of morphisms $f_i : A_i \to B_i$ such
that all
the diagrams
$$
\xymatrix{
A_i \ar[r]_{d_i} \ar[d]_{f_i} & A_{i - 1}
\ar[d]^{f_{i - 1}} \\
B_i \ar[r]^{d_i} & B_{i - 1}
}
$$
commute. The {\it category of chain complexes
of $\mathcal{A}$}
is denoted $\text{Ch}(\mathcal{A})$. The full
subcategory consisting
of objects of the form
$$
\ldots \to A_2 \to A_1 \to A_0 \to 0 \to 0
\to \ldots
$$
is denoted $\text{Ch}_{\geq 0}(\mathcal{A})$.
In other words, a chain complex $A_\bullet$
belongs to
$\text{Ch}_{\geq 0}(\mathcal{A})$ if and only
if
$A_i = 0$ for all $i < 0$.

\medskip\noindent
Given an additive category $\mathcal{A}$ we
identify $\mathcal{A}$
with the full subcategory of
$\text{Ch}(\mathcal{A})$ consisting
of chain complexes zero except in degree $0$
by the functor
$$
\mathcal{A} \longrightarrow
\text{Ch}(\mathcal{A}),\quad
A \longmapsto (\ldots \to 0 \to A \to 0 \to
\ldots)
$$
By abuse of notation we often denote the
object on the right hand side
simply $A$. If we want to stress that we are
viewing $A$ as a chain
complex we may sometimes use the notation
$A[0]$, see
Section \ref{section-homotopy-shift}.

\medskip\noindent
A {\it homotopy $h$} between a pair of
morphisms
of chain complexes $f, g : A_\bullet \to
B_\bullet$
is a collection of morphisms $h_i : A_i \to
B_{i + 1}$
such that we have
$$
f_i - g_i = d_{i + 1} \circ h_i + h_{i - 1}
\circ d_i
$$
for all $i$.
Two morphisms $f, g : A_\bullet \to
B_\bullet$ are
said to be {\it homotopic} if a homotopy
between $f$
and $g$ exists.
Clearly, the notions of chain complex,
morphism of
chain complexes, and homotopies between
morphisms of chain complexes
make sense even in a preadditive category.

\begin{lemma}
\label{lemma-compose-homotopy}
\begin{slogan}
Hom functors of $\text{Ch}(\mathcal{A})$
respect the homotopy relation.
\end{slogan}
Let $\mathcal{A}$ be an additive category.
Let $f, g : B_\bullet \to C_\bullet$ be
morphisms
of chain complexes. Suppose given morphisms
of chain
complexes $a : A_\bullet \to B_\bullet$, and
$c : C_\bullet \to D_\bullet$.
If $\{h_i : B_i \to C_{i + 1}\}$ defines a
homotopy
between $f$ and $g$, then $\{c_{i + 1} \circ
h_i \circ a_i\}$
defines a homotopy between $c \circ f \circ
a$ and
$c \circ g \circ a$.
\end{lemma}

\begin{proof}
Omitted.
\end{proof}

\noindent
In particular this means that it makes sense
to define
the category of chain complexes with maps up
to homotopy.
We'll return to this later.

\begin{definition}
\label{definition-homotopy-equivalent}
Let $\mathcal{A}$ be an additive category.
We say a morphism $a : A_\bullet \to
B_\bullet$
is a {\it homotopy equivalence} if there
exists
a morphism $b : B_\bullet \to A_\bullet$
such that there exists a homotopy between
$a \circ b$ and $\text{id}_A$
and there exists a homotopy between $b \circ
a$ and $\text{id}_B$.
If there exists such a morphism between
$A_\bullet$ and $B_\bullet$, then
we say that $A_\bullet$ and $B_\bullet$ are
{\it homotopy equivalent}.
\end{definition}

\noindent
In other words, two complexes are homotopy
equivalent if they become
isomorphic in the category of complexes up to
homotopy.

\begin{lemma}
\label{lemma-cat-chain-abelian}
Let $\mathcal{A}$ be an abelian category.
\begin{enumerate}
\item The category of chain complexes in
$\mathcal{A}$ is
abelian.
\item A morphism of complexes
$f : A_\bullet \to B_\bullet$ is injective
if and only if each $f_n : A_n \to B_n$ is
injective.
\item A morphism of complexes
$f : A_\bullet \to B_\bullet$ is surjective
if and only if each $f_n : A_n \to B_n$ is
surjective.
\item A sequence of chain complexes
$$
A_\bullet \xrightarrow{f} B_\bullet
\xrightarrow{g} C_\bullet
$$
is exact at $B_\bullet$ if and only if each
sequence
$$
A_i \xrightarrow{f_i} B_i \xrightarrow{g_i}
C_i
$$
is exact at $B_i$.
\end{enumerate}
\end{lemma}

\begin{proof}
Omitted.
\end{proof}

\noindent
For any $i \in \mathbf{Z}$ the $i$th {\it
homology group}
of a chain complex $A_\bullet$ in an abelian
category is defined by
the following formula
$$
H_i(A_\bullet) = \Ker(d_i)/\Im(d_{i + 1}).
$$
If $f : A_\bullet \to B_\bullet$ is a
morphism of chain
complexes of $\mathcal{A}$ then we get an
induced
morphism $H_i(f) : H_i(A_\bullet) \to
H_i(B_\bullet)$
because clearly
$f_i(\Ker(d_i : A_i \to A_{i - 1})) \subset
\Ker(d_i : B_i \to B_{i - 1})$, and similarly
for $\Im(d_{i + 1})$.
Thus we obtain a functor
$$
H_i : \text{Ch}(\mathcal{A}) \longrightarrow
\mathcal{A}.
$$

\begin{definition}
\label{definition-quasi-isomorphism}
Let $\mathcal{A}$ be an abelian category.
\begin{enumerate}
\item A morphism of chain complexes $f :
A_\bullet \to B_\bullet$
is called a {\it quasi-isomorphism} if the
induced
map $H_i(f) : H_i(A_\bullet) \to
H_i(B_\bullet)$
is an isomorphism for all $i \in \mathbf{Z}$.
\item A chain complex $A_\bullet$ is called
{\it acyclic} if all of its homology objects
$H_i(A_\bullet)$ are zero.
\end{enumerate}
\end{definition}


\begin{lemma}
\label{lemma-map-homology-homotopy}
Let $\mathcal{A}$ be an abelian category.
\begin{enumerate}
\item If the maps $f, g : A_\bullet \to
B_\bullet$ are
homotopic, then the induced maps $H_i(f)$ and
$H_i(g)$
are equal.
\item If the map $f : A_\bullet \to
B_\bullet$ is a homotopy
equivalence, then $f$ is a quasi-isomorphism.
\end{enumerate}
\end{lemma}

\begin{proof}
Omitted.
\end{proof}

\begin{lemma}
\label{lemma-long-exact-sequence-chain}
Let $\mathcal{A}$ be an abelian category.
Suppose that
$$
0 \to
A_\bullet \to
B_\bullet \to
C_\bullet \to
0
$$
is a short exact sequence of chain complexes
of $\mathcal{A}$.
Then there is a canonical long exact homology
sequence
$$
\xymatrix{
\ldots & \ldots & \ldots \ar[lld] \\
H_i(A_\bullet) \ar[r] & H_i(B_\bullet) \ar[r]
& H_i(C_\bullet) \ar[lld] \\
H_{i - 1}(A_\bullet) \ar[r] &
H_{i - 1}(B_\bullet) \ar[r] &
H_{i - 1}(C_\bullet) \ar[lld] \\
\ldots & \ldots & \ldots \\
}
$$
\end{lemma}

\begin{proof}
Omitted. The maps come from the Snake Lemma
\ref{lemma-snake}
applied to the diagrams
$$
\xymatrix{
&
A_i/\Im(d_{A, i + 1}) \ar[r] \ar[d]^{d_{A,
i}} &
B_i/\Im(d_{B, i + 1}) \ar[r] \ar[d]^{d_{B,
i}} &
C_i/\Im(d_{C, i + 1}) \ar[r] \ar[d]^{d_{C,
i}} &
0 \\
0 \ar[r] &
\Ker(d_{A, i - 1}) \ar[r] &
\Ker(d_{B, i - 1}) \ar[r] &
\Ker(d_{C, i - 1}) &
}
$$
\end{proof}

\begin{remark}
\label{remark-snake-lemma-les}
The snake lemma gives us the following:
let 
$$
\xymatrix{
0\ar[r]
&M'\ar[r]
&M\ar[r]
&M''\ar[r]
&0\\
& (P')^\bullet\ar[u]
&& (P'')^\bullet \ar[u]
}
$$
be a short exact sequence of $R$-modules
with $(P')^\bullet$ and
$(P'')^\bullet$ projective
resolutions. Then there is a
projective resolution $P^\bullet \to M$
such that
$$
\xymatrix{
0\ar[r]
&(P')^\bullet\ar[r]
&P^\bullet\ar[r]
&(P'')^\bullet\ar[r]
&0
}
$$
is exact in the corresponding category.
\end{remark}

\noindent
A {\it cochain complex $A^\bullet$} in an
additive category $\mathcal{A}$
is a complex
$$
\ldots \to
A^{n - 1} \xrightarrow{d^{n - 1}}
A^n \xrightarrow{d^n}
A^{n + 1} \to
\ldots
$$
of $\mathcal{A}$. In other words, we are
given an object $A^i$ of
$\mathcal{A}$ for all $i \in \mathbf{Z}$ and
for
all $i \in \mathbf{Z}$ a morphism $d^i : A^i
\to A^{i + 1}$ such that
$d^{i + 1} \circ d^i = 0$ for all $i$. A {\it
morphism of cochain
complexes $f : A^\bullet \to B^\bullet$} is
given by a
family of morphisms $f^i : A^i \to B^i$ such
that all
the diagrams
$$
\xymatrix{
A^i \ar[r]_{d^i} \ar[d]_{f^i} & A^{i + 1}
\ar[d]^{f^{i + 1}} \\
B^i \ar[r]^{d^i} & B^{i + 1}
}
$$
commute. The {\it category of cochain
complexes of $\mathcal{A}$}
is denoted $\text{CoCh}(\mathcal{A})$. The
full subcategory consisting
of objects of the form
$$
\ldots \to 0 \to 0 \to A^0 \to A^1 \to A^2
\to \ldots
$$
is denoted $\text{CoCh}_{\geq
0}(\mathcal{A})$.
In other words, a cochain complex $A^\bullet$
belongs to the subcategory
$\text{CoCh}_{\geq 0}(\mathcal{A})$ if and
only if
$A^i = 0$ for all $i < 0$.

\medskip\noindent
Given an additive category $\mathcal{A}$ we
identify $\mathcal{A}$
with the full subcategory of
$\text{CoCh}(\mathcal{A})$ consisting
of cochain complexes zero except in degree
$0$ by the functor
$$
\mathcal{A} \longrightarrow
\text{CoCh}(\mathcal{A}),\quad
A \longmapsto (\ldots \to 0 \to A \to 0 \to
\ldots)
$$
By abuse of notation we often denote the
object on the right hand side
simply $A$. If we want to stress that we are
viewing $A$ as a cochain
complex we may sometimes use the notation
$A[0]$, see
Section \ref{section-homotopy-shift}.

\medskip\noindent
A {\it homotopy $h$} between a pair of
morphisms
of cochain complexes $f, g : A^\bullet \to
B^\bullet$
is a collection of morphisms $h^i : A^i \to
B^{i - 1}$
such that we have
$$
f^i - g^i = d^{i - 1} \circ h^i + h^{i + 1}
\circ d^i
$$
for all $i$.
Two morphisms $f, g : A^\bullet \to
B^\bullet$ are
said to be {\it homotopic} if a homotopy
between $f$
and $g$ exists.
Clearly, the notions of cochain complex,
morphism of
cochain complexes, and homotopies between
morphisms of cochain complexes
make sense even in a preadditive category.

\begin{lemma}
\label{lemma-compose-homotopy-cochain}
Let $\mathcal{A}$ be an additive category.
Let $f, g : B^\bullet \to C^\bullet$ be
morphisms
of cochain complexes. Suppose given morphisms
of cochain
complexes $a : A^\bullet \to B^\bullet$, and
$c : C^\bullet \to D^\bullet$.
If $\{h^i : B^i \to C^{i - 1}\}$ defines a
homotopy
between $f$ and $g$, then $\{c^{i - 1} \circ
h^i \circ a^i\}$
defines a homotopy between $c \circ f \circ
a$ and
$c \circ g \circ a$.
\end{lemma}

\begin{proof}
Omitted.
\end{proof}

\noindent
In particular this means that it makes sense
to define
the category of cochain complexes with maps
up to homotopy.
We'll return to this later.

\begin{definition}
\label{definition-homot-equiv-cocha}
Let $\mathcal{A}$ be an additive category.
We say a morphism $a : A^\bullet \to
B^\bullet$
is a {\it homotopy equivalence} if there
exists
a morphism $b : B^\bullet \to A^\bullet$
such that there exists a homotopy between
$a \circ b$ and $\text{id}_A$
and there exists a homotopy between $b \circ
a$ and $\text{id}_B$.
If there exists such a morphism between
$A^\bullet$ and $B^\bullet$, then
we say that $A^\bullet$ and $B^\bullet$ are
{\it homotopy equivalent}.
\end{definition}

\noindent
In other words, two complexes are homotopy
equivalent if they become
isomorphic in the category of complexes up to
homotopy.

\begin{lemma}
\label{lemma-cat-cochain-abelian}
Let $\mathcal{A}$ be an abelian category.
\begin{enumerate}
\item The category of cochain complexes in
$\mathcal{A}$ is
abelian.
\item A morphism of cochain complexes
$f : A^\bullet \to B^\bullet$ is injective
if and only if each $f^n : A^n \to B^n$ is
injective.
\item A morphism of cochain complexes
$f : A^\bullet \to B^\bullet$ is surjective
if and only if each $f^n : A^n \to B^n$ is
surjective.
\item A sequence of cochain complexes
$$
A^\bullet \xrightarrow{f} B^\bullet
\xrightarrow{g} C^\bullet
$$
is exact at $B^\bullet$ if and only if each
sequence
$$
A^i \xrightarrow{f^i} B^i \xrightarrow{g^i}
C^i
$$
is exact at $B^i$.
\end{enumerate}
\end{lemma}

\begin{proof}
Omitted.
\end{proof}

\noindent
For any $i \in \mathbf{Z}$ the $i$th {\it
cohomology group}
of a cochain complex $A^\bullet$ is defined
by
the following formula
$$
H^i(A^\bullet) = \Ker(d^i)/\Im(d^{i - 1}).
$$
If $f : A^\bullet \to B^\bullet$ is a
morphism of cochain
complexes of $\mathcal{A}$ then we get an
induced
morphism $H^i(f) : H^i(A^\bullet) \to
H^i(B^\bullet)$
because clearly
$f^i(\Ker(d^i : A^i \to A^{i + 1})) \subset
\Ker(d^i : B^i \to B^{i + 1})$, and similarly
for $\Im(d^{i - 1})$.
Thus we obtain a functor
$$
H^i : \text{CoCh}(\mathcal{A})
\longrightarrow \mathcal{A}.
$$

\begin{definition}
\label{definition-quasi-isomorphism-cochain}
Let $\mathcal{A}$ be an abelian category.
\begin{enumerate}
\item A morphism of cochain complexes $f :
A^\bullet \to B^\bullet$
of $\mathcal{A}$ is called a {\it
quasi-isomorphism} if the induced
maps $H^i(f) : H^i(A^\bullet) \to
H^i(B^\bullet)$
is an isomorphism for all $i \in \mathbf{Z}$.
\item A cochain complex $A^\bullet$ is called
{\it acyclic} if all of its cohomology
objects
$H^i(A^\bullet)$ are zero.
\end{enumerate}
\end{definition}

\begin{lemma}
\label{lemma-map-cohomology-homotopy-cochain}
Let $\mathcal{A}$ be an abelian category.
\begin{enumerate}
\item If the maps $f, g : A^\bullet \to
B^\bullet$ are
homotopic, then the induced maps $H^i(f)$ and
$H^i(g)$
are equal.
\item If $f : A^\bullet \to B^\bullet$ is a
homotopy equivalence,
then $f$ is a quasi-isomorphism.
\end{enumerate}
\end{lemma}

\begin{proof}
Omitted.
\end{proof}

\begin{lemma}
\label{lemma-long-exact-sequence-cochain}
\begin{slogan}
Short exact sequences of complexes give rise
to long exact sequences
of (co)homology.
\end{slogan}
Let $\mathcal{A}$ be an abelian category.
Suppose that
$$
0 \to
A^\bullet \to
B^\bullet \to
C^\bullet \to
0
$$
is a short exact sequence of cochain
complexes of $\mathcal{A}$.
Then there is a long exact cohomology
sequence
$$
\xymatrix{
\ldots & \ldots & \ldots \ar[lld] \\
H^i(A^\bullet) \ar[r] &
H^i(B^\bullet) \ar[r] &
H^i(C^\bullet) \ar[lld] \\
H^{i + 1}(A^\bullet) \ar[r] &
H^{i + 1}(B^\bullet) \ar[r] &
H^{i + 1}(C^\bullet) \ar[lld] \\
\ldots & \ldots & \ldots \\
}
$$
The construction produces long exact
cohomology sequences
which are functorial in the short exact
sequence and compatible with shifts as in
Definition \ref{definition-cohomology-shift}.
\end{lemma}

\begin{proof}
For the horizontal maps $H^i(A^\bullet) \to
H^i(B^\bullet)$ and
$H^i(B^\bullet) \to H^i(C^\bullet)$ we use
the fact that $H^i$
is a functor, see above. For the ``boundary
map''
$H^i(C^\bullet) \to H^{i + 1}(A^\bullet)$ we
use the map $\delta$
of the Snake Lemma \ref{lemma-snake}
applied to the diagram
$$
\xymatrix{
&
A^i/\Im(d_A^{i - 1}) \ar[r] \ar[d]^{d_A^i} &
B^i/\Im(d_B^{i - 1}) \ar[r] \ar[d]^{d_B^i} &
C^i/\Im(d_C^{i - 1}) \ar[r] \ar[d]^{d_C^i} &
0 \\
0 \ar[r] &
\Ker(d_A^{i + 1}) \ar[r] &
\Ker(d_B^{i + 1}) \ar[r] &
\Ker(d_C^{i + 1}) &
}
$$
This works as the kernel of the right
vertical map is equal to
$H^i(C^\bullet)$ and the cokernel of the left
vertical map is
$H^{i + 1}(A^\bullet)$. The exactness of the
long sequence is
the exactnesss in part (2) of Lemma
\ref{lemma-snake}.
The functoriality is Lemma
\ref{lemma-snake-natural}. Compatibility
with shifts is immediate from the
definitions.
\end{proof}


\section{Injectives}
\label{section-injectives}

\noindent
The upshot of injectives is
that whenever there is a map
defined on a subobject $A \subset B$,
then we can extend the map to the whole $B$.
Note that this defining property is not
universal: the extension may not be unique.

Also note that we discussed this concept,
and also the projective concept,
in Commutative Algebra.

\begin{definition}
\label{definition-injective}
Let $\mathcal{A}$ be an abelian category.
An object $J \in \Ob(\mathcal{A})$ is
called {\it injective} if for every
injection $A \hookrightarrow B$ and every
morphism $A \to J$ there exists a morphism
$B \to J$ making the following diagram
commute:
$$
\xymatrix{
A \ar[r] \ar[d] &
B \ar@{-->}[ld] \\
J &
}
$$
\end{definition}

\noindent
Here is the obligatory characterization of
injective objects.

\begin{lemma}
\label{lemma-characterize-injectives}
Let $\mathcal{A}$ be an abelian category,
and let $I$ be an object of $\mathcal{A}$.
The following are equivalent:
\begin{enumerate}
\item The object $I$ is injective.
\item The functor
$B \mapsto \Hom_\mathcal{A}(B,I)$ is exact.
\item Every short exact sequence
$$
0 \to I \to A \to B \to 0
$$
in $\mathcal{A}$ is split.
\item For every
$B \in \Ob(\mathcal{A})$, we have
$$
\Ext_\mathcal{A}(B,I)=0.
$$
\end{enumerate}
\end{lemma}

\begin{proof}
Omitted.
\end{proof}

\begin{lemma}
\label{lemma-product-injectives}
Let $\mathcal{A}$ be an abelian category.
Suppose that $I_\omega$, for
$\omega \in \Omega$, is a set of injective
objects of $\mathcal{A}$. If
$$
\prod_{\omega \in \Omega} I_\omega
$$
exists, then it is injective.
\end{lemma}

\begin{proof}
Omitted.
\end{proof}

\begin{definition}
\label{definition-enough-injectives}
Let $\mathcal{A}$ be an abelian category.
We say that $\mathcal{A}$ has {\it enough
injectives} if every object $A$ has an
injective morphism $A \to J$ into an
injective object $J$.
\end{definition}

\begin{definition}
\label{definition-injective-resolution}
Let $\mathcal{A}$ be an abelian category.
Let $A \in \Ob(\mathcal{A})$.
An {\it injective resolution of $A$} is a
complex $I^\bullet$ together with a map
$A \to I^0$ such that:
\begin{enumerate}
\item We have $I^n = 0$ for $n < 0$.
\item Each $I^n$ is an injective object of
$\mathcal{A}$.
\item The map $A \to I^0$ is an isomorphism
onto $\Ker(d^0)$.
\item We have $H^i(I^\bullet) = 0$ for
$i > 0$.
\end{enumerate}
Hence $A[0] \to I^\bullet$ is a
quasi-isomorphism.
In other words the complex
$$
\ldots \to 0 \to A \to I^0 \to I^1
\to \ldots
$$
is acyclic.
\end{definition}

\begin{lemma}
\label{lemma-injective-resolutions-exist}
Let $\mathcal{A}$ be an abelian category.
Assume $\mathcal{A}$ has enough injectives.
\begin{enumerate}
\item Any object of $\mathcal{A}$ has an
injective resolution.
\end{enumerate}
\end{lemma}

\begin{proof}
Proof of (1). First choose an injection
$A \to I^0$ of $A$ into an injective object
of $\mathcal{A}$. Next, choose an injection
$I_0/A \to I^1$ into an injective object of
$\mathcal{A}$. Denote $d^0$ the induced map
$I^0 \to I^1$.
Next, choose an injection
$I^1/\Im(d^0) \to I^2$ into an injective
object of $\mathcal{A}$. Denote $d^1$ the
induced map $I^1 \to I^2$. And so on.
\end{proof}

\begin{definition}
\label{definition-functorial-injectives}
Let $\mathcal{A}$ be an abelian category.
We say that $\mathcal{A}$ has {\it
functorial injective embeddings} if there
exists a functor $J$ from $\mathcal{A}$ to
its arrow category such that:
\begin{enumerate}
\item the source of $J(A)$ is $A$;
\item for every object
$A \in \Ob(\mathcal{A})$, the morphism
$J(A)$ is injective;
\item for every object
$A \in \Ob(\mathcal{A})$, the target of
$J(A)$ is an injective object of
$\mathcal{A}$.
\end{enumerate}
We denote such a functor by
$A \mapsto (A \to J(A))$.
\end{definition}








\section{Derived functors}
\label{section-derived-functos}

\noindent
Basic sort of question in algebra is,
say some object $C=B/A$ defined
as a quotient,
and we have some process $F$
(a functor), does $F(C)=F(B)/F(A)$,
and if not, how to correct this naive guess.

More abstractly, if
$$
\xymatrix{
0\ar[r]
&A\ar[r]^i
&B\ar[r]^p
&C\ar[r]
&0
}
$$
is a short exact sequence, then does
$$
\xymatrix{
0\ar[r]
&F(A)\ar[r]
&F(B)\ar[r]
&F(C)\ar[r]
&0
}
$$

\begin{remark}
\label{remark-exact-functors}
In homological algebra,
exact functors are kind of rare,
but it is very common that a functor be
either left exact or right exact.
\end{remark}

\begin{definition}
\label{definition-exact-functors}
Let $\mathcal{A}$ and $\mathcal{B}$ be
abelian categories, and let
$F : \mathcal{A} \to \mathcal{B}$ be a
functor. For every short exact sequence
$$
0 \to A \to B \to C \to 0
$$
in $\mathcal{A}$, consider the induced
sequence in $\mathcal{B}$.
\begin{enumerate}
\item The functor $F$ is {\it left exact}
if
$$
0 \to F(A) \to F(B) \to F(C)
$$
is exact.
\item The functor $F$ is {\it right exact}
if
$$
F(A) \to F(B) \to F(C) \to 0
$$
is exact.
\item The functor $F$ is {\it exact} if
$$
0 \to F(A) \to F(B) \to F(C) \to 0
$$
is exact.
\end{enumerate}
\end{definition}

\begin{definition}
\label{definition-right-derived-functors}
Let $\mathcal{A}$ have enough injectives.
Let $F: \mathcal{A} to \mathcal{B}$ be a
left exact functor.
Define the {\it right derived functors}
of $F$ to be
$$
R^nF(A) = H^n(F(I^\bullet)),
$$
when $0 \to I^0 \to I^1 \to I^2 \to…$
is an injective resolution of $\mathcal{A}$.
\end{definition}

\begin{theorem}
\label{theorem-derived-functors}
\begin{enumerate}
\item The value of $R^nF(A) \in \mathcal{B}$
does not depend on the choice
of $A \to I^\bullet$ (up to isomorphism).

\item A morphism $f:A \to A'$
in $\mathcal{A}$ induces
$R^nF(A) \to R^nF(A')$
well-defined, and $R^nF$,
$n \geq  0$ are all functors
$\mathcal{A} \to \mathcal{B}$.

\item $R^0F(A) = F(A)$.

\item If
$0 \to A \to B \to C \to 0$
is a short exact sequence, then we have

$$
\xymatrix{
&R^2F(A) \ar[r] & \cdots\\
&R^1F(A) \ar[r] & R^1F(B)\ar[r]&R^1F(C)
\ar[ull]\\
0\ar[r] &
F(A)\ar[r]&F(B)\ar[r]&F(C)\ar[ull]^\delta
}
$$
is exact.
\end{enumerate}
\end{theorem}


\medskip\noindent
{\bf Divisible groups.}

\begin{definition}
\label{definition-divisible-group}
An abelian group $I$ is called
{\it divisible} if $\forall i \in I$
and $n \in \mathbb{Z}\setminus\{0\}$
there exists $i' \in I$ such that
$ni'=i$.
\end{definition}

\noindent
Example: $I = \mathbb{Q}$.

\section{Projectives}
\label{section-projectives}

\noindent
Projective objects have a lifting property
dual to the extension property of injective
objects. A morphism from a projective
object to a quotient lifts to the object
being quotiented. The lift need not be
unique.

\begin{definition}
\label{definition-projective}
Let $\mathcal{A}$ be an abelian category.
An object $P$ of $\mathcal{A}$ is called
{\it projective} if, for every epimorphism
$f:A\to B$ and every morphism
$\alpha:P\to B$, there is a morphism
$\beta:P\to A$ such that
$$
\alpha=f\circ\beta.
$$
Thus there is a commutative diagram
$$
\xymatrix{
A \ar[r]^f & B \ar[r] & 0 \\
P \ar[u]^{\beta} \ar[ur]_{\alpha} &
}
$$
for every such $f$ and $\alpha$.
\end{definition}

\begin{definition}
\label{definition-projection-module}
Let $R$ be a ring and let $F$ be an
$R$-module. A {\it projection on $F$} is
an endomorphism $\pi:F\to F$ satisfying
$$
\pi^2=\pi.
$$
\end{definition}

\begin{lemma}
\label{lemma-projection-decomposition}
Let $\pi:F\to F$ be a projection on an
$R$-module. Then
$$
F\cong\Ker(\pi)\oplus\Im(\pi).
$$
\end{lemma}

\begin{proof}
Set $K=\Ker(\pi)$ and $I=\Im(\pi)$.
If $x\in K\cap I$, then $x=\pi(y)$ for
some $y\in F$, and hence
$$
x=\pi(y)=\pi^2(y)=\pi(x)=0.
$$
On the other hand, every $x\in F$ has the
decomposition
$$
x=\pi(x)+(x-\pi(x)),
$$
where the first summand belongs to $I$ and
the second belongs to $K$.
\end{proof}

\begin{definition}
\label{definition-projective-direct-summand}
Let $R$ be a ring. A {\it projective
$R$-module} is an $R$-module $P$ for which
there are an $R$-module $C$, a free
$R$-module $F$, and an isomorphism
$$
F\cong P\oplus C.
$$
\end{definition}

\begin{lemma}[Corollary]
\label{lemma-projection-projective}
The kernel and image of a projection on a
free $R$-module are projective
$R$-modules.
\end{lemma}

\begin{proof}
This follows from Lemma
\ref{lemma-projection-decomposition}.
\end{proof}

\begin{example}
\label{example-moebius-projective}
Consider the ring
$$
R=
\mathbb{R}[x,y]/(x^2+y^2-1)
$$
and the endomorphism of $R^2$ given by
$$
\pi
=
\frac{1}{2}
\begin{pmatrix}
1+x & y \\
y & 1-x
\end{pmatrix}.
$$
The relation in $R$ gives $\pi^2=\pi$.
Consequently, $\Im(\pi)$ is a projective
$R$-module.

Geometrically, $R$ is the ring of
polynomial functions on the real unit
circle. At the point
$$
p(\theta)
=
(\cos(\theta),\sin(\theta)),
$$
the image of $\pi$ is the line spanned by
$$
(\cos(\theta/2),\sin(\theta/2)).
$$
As the point goes once around the circle,
this line makes a half turn. Thus the
module describes the sections of the
Möbius line bundle.
\end{example}

\begin{lemma}
\label{lemma-projective-module-object}
An object of $\mathit{Mod}_R$ is projective
in the sense of Definition
\ref{definition-projective} if and only if
it is a projective $R$-module in the sense
of Definition
\ref{definition-projective-direct-summand}.
\end{lemma}

\begin{proof}
Suppose first that $P$ is projective as an
object of $\mathit{Mod}_R$. Choose a free
$R$-module $F$ with an epimorphism
$f:F\to P$. The identity of $P$ lifts to
a morphism $s:P\to F$ satisfying
$$
f\circ s=\operatorname{id}_P.
$$
In particular, $s$ is injective. If
$K=\Ker(f)$, every $x\in F$ can be written
uniquely as
$$
x=(x-s(f(x)))+s(f(x)),
$$
with the first summand in $K$ and the
second in $s(P)$. Moreover,
$K\cap s(P)=0$. Thus
$$
F\cong K\oplus P.
$$

Conversely, suppose that
$F\cong P\oplus K$ for a free module $F$.
Let $f:A\to B$ be an epimorphism and let
$\alpha:P\to B$ be a morphism. Extend
$\alpha$ to a morphism $F\to B$ by making
it zero on $K$. A map from a free module
lifts across every epimorphism: choose a
basis of $F$ and a preimage in $A$ of the
image of each basis element. Restricting
such a lift $F\to A$ to $P$ gives the
required lift of $\alpha$.
\end{proof}

\begin{definition}
\label{definition-enough-projectives}
An abelian category $\mathcal{A}$ has
{\it enough projectives} if every object
$A$ admits an epimorphism $P\to A$ from a
projective object $P$.
\end{definition}

\begin{lemma}
\label{lemma-modules-enough-projectives}
For every ring $R$, the category
$\mathit{Mod}_R$ has enough projectives.
\end{lemma}

\begin{proof}
Every $R$-module is the quotient of a free
$R$-module. Every free $R$-module is
projective by Lemma
\ref{lemma-projective-module-object}.
\end{proof}

\begin{definition}
\label{definition-projective-resolution}
Let $\mathcal{A}$ be an abelian category
and let $A$ be an object of $\mathcal{A}$.
A {\it projective resolution of $A$} is a
chain complex $P_\bullet$ together with an
augmentation $P_0\to A$ such that:
\begin{enumerate}
\item We have $P_n=0$ for $n<0$.
\item Every $P_n$ is projective.
\item The augmented complex
$$
\ldots\to P_2\to P_1\to P_0
\to A\to 0
$$
is exact.
\end{enumerate}
Equivalently, $P_\bullet\to A[0]$ is a
quasi-isomorphism.
\end{definition}

\begin{lemma}
\label{lemma-projective-resolutions-exist}
If an abelian category $\mathcal{A}$ has
enough projectives, then every object of
$\mathcal{A}$ has a projective resolution.
In particular, every $R$-module has a
projective resolution.
\end{lemma}

\begin{proof}
Choose an epimorphism $P_0\to A$ with
$P_0$ projective. Next choose an
epimorphism
$$
P_1\to\Ker(P_0\to A)
$$
with $P_1$ projective, and compose it with
the inclusion into $P_0$. Continuing with
the kernel at each stage produces an exact
augmented complex
$$
\ldots\to P_2\to P_1\to P_0
\to A\to 0.
$$
The last assertion follows from Lemma
\ref{lemma-modules-enough-projectives}.
\end{proof}

\noindent
Projective resolutions are the dual input
to injective resolutions. If a right exact
functor admits left derived functors, they
are computed by applying the functor to a
projective resolution and taking homology.

\section{Group cohomology}
\label{section-group-cohomology}

\noindent
We have considered the category
$G$-mod $\sim$ $\mathbb{Z}G$-mod and the
functor of $G$-invariants
$$
(-)^G : G\operatorname{-mod}
\to
\underline{\operatorname{Ab}}.
$$
this functor is left exact, and
$\mathbb{Z}G$-mod has enough injectives,
so we can construct right derived functors
$$
R^n(-)^G: G\operatorname{-mod}
\to \underline{\operatorname{Ab}}.
$$
More generally, if $R$ is a ring and
$M$ an $R$-module,
$\Hom_R(M,-):R\operatorname{-mod}
\to \underline{ \operatorname{Ab}}$,
and 
$[R^n\Hom_R(M,-)](A)$
is called $\operatorname{Ext}^n_R(M,A)$.

\begin{remark}
\label{remark-hom-functor-contravariant}
Fixing an $R$-module $N$,
the functor
$$
\Hom_R(-,N):R\operatorname{-mod}
\to \underline{\operatorname{Ab}}
$$
is a left exact (and contravariant)
functor.
\end{remark}

\noindent
So if $R$-mod has enough projectives
(which it does) we can define left
derived functors
$$
[L^n\Hom_R(-,N)](M).
$$

\begin{theorem}[Balancing]
\label{theorem-balancing}
$$
[R^n\Hom_R(M,-)](N)
\simeq
[L^n\Hom_R(-,N)](M)
$$
and both equal
$$
\operatorname{Ext}^n_R(M,N).
$$
\end{theorem}

\noindent
We shall prove this theorem later.

\medskip\noindent
Back to group cohomology,
$$
R^n(-)^G](M)
=
\operatorname{Ext}^n_{\mathbb{Z} G}
(\mathbb{Z},M).
$$
We can compute it by forming an injective
resolution of $M$… OR (!)
a projective resolution of $\mathbb{Z}$
(which is better).



\medskip\noindent
{\bf The bar complex.}
The projective resolution of $\mathbb{Z}$
will be
$$
… \xrightarrow{d}P_2
\xrightarrow{d}P_1
\xrightarrow{d}P_0
\xrightarrow{\varepsilon}
\mathbb{Z}\to 0
$$
with $P_n=\mathbb{Z}G^{n+1}$
as abelian group (free $\mathbb{Z}$-module).

$G$-action on $\mathbb{Z}G^{n+1}$ is
$$
g(g_0,g_1,…,g_n)
=
(gg_0, g g_1,…,g g_n),
$$
extended $\mathbb{Z}$-linearly.

The differential
$$
\mathbb{Z}G^{n+1}
\xrightarrow{d}
\mathbb{Z}G^n,
$$
$$
d((g_0,g_1,…,g_n))
=
\sum_{i=0}^n
(-1)^i(g_0,g_1,…,\widehat{g_i},…,g_n).
$$
So for example
$$
\begin{aligned}
\mathbb{Z}G^2 &\longrightarrow \mathbb{Z}G \\
(g_0,g_1) &\longmapsto g_1-g_0,
\end{aligned}
$$
and
$$
\begin{aligned}
\mathbb{Z}G^3 &\longrightarrow \mathbb{Z}G \\
(g_0,g_1,g_2) &\longmapsto 
(g_1,g_2)-(g_0,g_2)+(g_0,g_1).
\end{aligned}
$$
The ``accumulation'' is
$$
\varepsilon:\mathbb{Z}G \to \mathbb{Z},
\qquad 
\varepsilon(g)=1 \quad \forall g \in G.
$$
\begin{exercise}
\label{exercise-bar-complex}
$d \circ d=0$.
\end{exercise}

\begin{remark}
\label{remark-bar}
Every $P_n$ is a free $\mathbb{Z}G$-module,
and free modules are projective.
Why is $\tilde{P}_\bullet$
acyclic? We'll see soon.
\end{remark}

\noindent
This leads to the following
definition/theorem.

\begin{definition}
\label{definition-group-cohomology}
Let $A$ be a $G$-module.
We define
$$
C^n(G,A)
=
\{\varphi:G^n\to A\}
$$
(no restrictions whatsoever, the $\varphi$
are just functions of sets).
$C^n$ is an abelian group (by the addition
of $A$). Let $C^n = C^n(G,A)$.
Differentials
$$
0\to C^0
\xrightarrow{d} C^1
\xrightarrow{d} C^2
\xrightarrow{d}…
$$
If $\varphi \in C^n$,
$$
[d\varphi](g_0,g_1,…,g_n)
=
g_0\varphi(g_1,g_2,…,g_n)
-
\sum_{i=a}^{n-1}
(-1)^i
\varphi(g_0|g_1|…g_{i-1}|g_ig_{i+1}|
g_{i+2}|…g_n)
+(-1)^{n+1}\varphi(g_0|g_1|…|g_{n-1}).
$$
Let $Z^n=\ker d$,
$B^n=\operatorname{Im}(d)$, since
$d^2=0$,  $B^n \subset Z^n$.
The group cohomology is
$$
H^n(G,A) = Z^n/B^n.
$$
\end{definition}

\begin{exercise}
\label{exercise-hom-group-cohom}
$$
\Hom_{\mathbb{Z}G}(P_n,A)
=
C^n(G,A).
$$
\end{exercise}

\begin{theorem}
\label{theorem-ext-group-cohom}
$$
\operatorname{Ext}^n_{\mathbb{Z}A}(\mathbb{Z},A)
=
H^n(G,A).
$$
\end{theorem}

\begin{proof}
\begin{enumerate}
\item Resolve $\mathbb{Z}$ by 
$P_{\bullet}$ (balancing says this is
legit).

\item Exercise says
$$
\operatorname{Ext}^n_{\mathbb{Z}G}(\mathbb{Z},A)
=
[L^n\Hom_{\mathbb{Z}G}(Z,-)](A)
=
\Hom_{\mathbb{Z}G}(P_n,A)
=
C^n(G,A).
$$
\item Remains to prove
$$
\tilde{P}_\bullet=P_\bullet
\to \mathbb{Z}
\xrightarrow{\varepsilon}0
$$
is acyclic. [Proof continued.]
\end{enumerate}
\end{proof}

\noindent
Unwrapping definitions we see that
$$
H^0(G,A) = A^G,
$$
$$
H^1(G,A)
=
\frac{\{\varphi:G \to A:
g\varphi(h)-\varphi(gh)
+\varphi(g)=0\}}{
\varphi \operatorname{da}
\operatorname{forma}
\varphi(g)=ga-a}.
$$
The bar resolution is very big, its
advantage is it works for all $G$.
If I'm intersete in a specific $G$,
maybe there exists a smaller resolution…

\begin{example}
\label{example-group-cohom-cyclic}
Let $G = C_m \simeq \mathbb{Z}/m$
be the cylcic group…
\end{example}


\noindent
Anything else in this lecture?


\section{Extensions}
\label{section-extensions}

\begin{definition}
\label{definition-extension}
Let $\mathcal{A}$ be an abelian category.
Let $A, B \in \Ob(\mathcal{A})$.
An {\it extension $E$ of $B$ by $A$} is a
short
exact sequence
$$
0 \to A \to E \to B \to 0.
$$
A {\it morphism of extensions} between two
extensions $0 \to A \to E \to B \to 0$ and
$0 \to A \to F \to B \to 0$ means a morphism
$f : E \to F$ in $\mathcal{A}$ making the
diagram
$$
\xymatrix{
0 \ar[r] &
A \ar[r] \ar[d]^{\text{id}} &
E \ar[r] \ar[d]^f &
B \ar[r] \ar[d]^{\text{id}} &
0 \\
0 \ar[r] &
A \ar[r] &
F \ar[r] &
B \ar[r] &
0
}
$$
commutative.
Thus, the extensions of $B$ by $A$ form a
category.
\end{definition}

\noindent
By abuse of language we often omit mention of
the
morphisms $A \to E$ and $E \to B$, although
they are
definitively part of the structure of an
extension.

\begin{definition}
\label{definition-ext-group}
Let $\mathcal{A}$ be an abelian category.
Let $A, B \in \Ob(\mathcal{A})$.
The set of isomorphism classes of extensions
of $B$ by $A$ is denoted
$$
\Ext_\mathcal{A}(B, A).
$$
This is called the {\it $\Ext$-group}.
\end{definition}

\noindent
This definition works, because by our
conventions
$\Ob(\mathcal{A})$ is a set, and hence
$\Ext_\mathcal{A}(B, A)$ is a set.
In any of the cases of ``big'' abelian
categories
listed in Categories, Remark
\ref{categories-remark-big-categories}
one can check by hand that
$\Ext_\mathcal{A}(B, A)$
is a set as well. Also, we will see later
that this is
always the case when $\mathcal{A}$ has either
enough projectives
or enough injectives. Insert future reference
here.

\medskip\noindent
Actually we can turn $\Ext_\mathcal{A}(-, -)$
into a
functor
$$
\mathcal{A} \times \mathcal{A}^{opp}
\longrightarrow \textit{Sets}, \quad
(A, B) \longmapsto \Ext_\mathcal{A}(B, A)
$$
as follows:
\begin{enumerate}
\item Given a morphism $B' \to B$ and an
extension
$E$ of $B$ by $A$ we define $E' = E \times_B
B'$.
By Lemma \ref{lemma-cartesian-kernel} we have
the following
commutative diagram of short exact sequences
$$
\xymatrix{
0 \ar[r] & A \ar[r] \ar[d] & E' \ar[r] \ar[d]
& B' \ar[r] \ar[d] & 0 \\
0 \ar[r] & A \ar[r] & E \ar[r] & B \ar[r] & 0
}
$$
The extension $E'$ is called the {\it
pullback of $E$ via
$B' \to B$}.
\item Given a morphism $A \to A'$ and an
extension
$E$ of $B$ by $A$ we define $E' = A' \amalg_A
E$.
By Lemma \ref{lemma-cartesian-cocartesian}
we have the following commutative diagram
of short exact sequences
$$
\xymatrix{
0 \ar[r] & A \ar[r] \ar[d] & E \ar[r] \ar[d]
& B \ar[r] \ar[d] & 0 \\
0 \ar[r] & A' \ar[r] & E' \ar[r] & B \ar[r] &
0
}
$$
The extension $E'$ is called the {\it pushout
of $E$ via
$A \to A'$}.
\end{enumerate}
To see that this defines a functor as
indicated above
there are several things to verify. First of
all
functoriality in the variable $B$ requires
that
$(E \times_B B') \times_{B'} B'' = E \times_B
B''$
which is a general property of fibre
products.
Dually one deals with functoriality in the
variable $A$. Finally, given $A \to A'$ and
$B' \to B$ we have to show that
$$
A' \amalg_A (E \times_B B')
\cong
(A' \amalg_A E)\times_B B'
$$
as extensions of $B'$ by $A'$. Recall that
$A' \amalg_A E$
is a quotient of $A' \oplus E$.
Thus the right hand side is a quotient of
$A' \oplus E \times_B B'$, and it is
straightforward to see that
the kernel is exactly what you need in order
to
get the left hand side.

\medskip\noindent
Note that if $E_1$ and $E_2$ are extensions
of
$B$ by $A$, then $E_1\oplus E_2$ is an
extension
of $B \oplus B$ by $A\oplus A$. We
push out by the sum map $A \oplus A \to A$
and we
pull back by the diagonal map $B \to B \oplus
B$ to get
an extension $E_1 + E_2$ of $B$ by $A$.
$$
\xymatrix{
0 \ar[r] &
A \oplus A \ar[r] \ar[d]_\Sigma &
E_1 \oplus E_2 \ar[r] \ar[d] &
B \oplus B \ar[r] \ar[d] &
0 \\
0 \ar[r] &
A \ar[r] &
E' \ar[r] &
B \oplus B \ar[r] &
0\\
0 \ar[r] &
A \ar[r] \ar[u] &
E_1 + E_2 \ar[r] \ar[u] &
B \ar[r] \ar[u]^\Delta &
0
}
$$
The extension $E_1 + E_2$ is called the {\it
Baer sum} of the
given extensions.

\begin{lemma}
\label{lemma-baer-sum}
The construction $(E_1, E_2) \mapsto E_1 +
E_2$
above defines a commutative group
law on $\Ext_\mathcal{A}(B, A)$ which is
functorial in both variables.
\end{lemma}

\begin{proof}
Omitted.
\end{proof}

\begin{lemma}
\label{lemma-six-term-sequence-ext}
Let $\mathcal{A}$ be an abelian category.
Let $0 \to M_1 \to M_2 \to M_3 \to 0$ be a
short exact sequence
in $\mathcal{A}$.
\begin{enumerate}
\item There is a canonical six term exact
sequence of abelian groups
$$
\xymatrix{
0 \ar[r] &
\Hom_\mathcal{A}(M_3, N) \ar[r] &
\Hom_\mathcal{A}(M_2, N) \ar[r] &
\Hom_\mathcal{A}(M_1, N) \ar[lld] \\
& \Ext_\mathcal{A}(M_3, N) \ar[r] &
\Ext_\mathcal{A}(M_2, N) \ar[r] &
\Ext_\mathcal{A}(M_1, N)
}
$$
for all objects $N$ of $\mathcal{A}$, and
\item there is a canonical six term exact
sequence of abelian groups
$$
\xymatrix{
0 \ar[r] &
\Hom_\mathcal{A}(N, M_1) \ar[r] &
\Hom_\mathcal{A}(N, M_2) \ar[r] &
\Hom_\mathcal{A}(N, M_3) \ar[lld] \\
& \Ext_\mathcal{A}(N, M_1) \ar[r] &
\Ext_\mathcal{A}(N, M_2) \ar[r] &
\Ext_\mathcal{A}(N, M_3)
}
$$
for all objects $N$ of $\mathcal{A}$.
\end{enumerate}
\end{lemma}

\begin{proof}
Omitted. Hint: The boundary maps are defined
using either the pushout
or pullback of the given short exact
sequence.
\end{proof}

\section{Example: Group cohomology 4-sept}
\label{section-group-cohomology-4-sept}

\noindent
What one would like to do is apply the LES of
group cohomology to deduce vanishing of
higher cohomology directly from exactness of
$(-)^G$.
Indeed if $V$ is a representation of $G$ one
could consider an embedding $V \to I$ into an
injective object in the category
$\operatorname{Rep}_F(G)$ and corresponding
quotient $Q=I/V$. The LES
$$
\xymatrix{
H^1(G,V) \ar[r] & H^1(G,I)=0 \\
0 \ar[r] & V^G \ar[r] & I^G \ar[r] & Q^G
\ar@{-->}[ulll]^0
}
$$
gives us the desired conclusion.

\noindent
This suggests two issues. Does
$\operatorname{Rep}(G)$
have enough injectives?
Do the derived functors of $(-)^G$ computed
within $\operatorname{Rep}(G)$ coincide with
those computed in $G$-mod? A better version
of
the second question is whether an object of
$\operatorname{Rep}(G)$ can be resolved by
injectives in $\operatorname{Rep}(G)$
which are
also injective in $G$-mod.

\noindent
This is a good excuse to introduce a bit of
abstract nonsense.

\begin{proposition}
\label{proposition-forgetful-4-sept}
The forgetful functor
$\operatorname{Rep}(G)\to\operatorname{Vec}$
has a right adjoint.
\end{proposition}

\begin{proof}
The right adjoint, explicitly, is the
following
so-called coinduction:
$$
\operatorname{Coind}(W)
=
\operatorname{Hom}_F(FG,W).
$$
The $G$-action is inherited from that on the
first factor. More precisely, if
$\varphi:FG\to W$, then $g\varphi$ is by
definition the function $FG\to W$ defined by
$$
x\longmapsto\varphi(xg)
$$
for all $x\in FG$.

\noindent
Let
$$
f\in
\operatorname{Hom}_{\operatorname{Rep}(G)}
(M,\operatorname{Hom}_F(FG,W)),
$$
and denote $f(m)$ by $\varphi_m$. Then by
evaluating $\varphi_m$ on $e\in FG$ we get an
element $\varphi_m(e)\in W$, and this gives
us
a function $\overline{f}:M\to W$. In this way
we have set up a function
$$
\operatorname{Hom}_{\operatorname{Rep}(G)}
(M,\operatorname{Hom}_F(FG,W))
\longrightarrow
\operatorname{Hom}_F(M,W)
$$
by $f\longmapsto\overline{f}$. Note that
$\overline{f}$, while linear in $m\in M$, is
not necessarily compatible with the
$G$-actions,
so it lies in $\operatorname{Hom}_F(M,W)$
rather than in
$\operatorname{Hom}_{\operatorname{Rep}(G)}
(M,W)$. Or even more precisely, it lies in
$\operatorname{Hom}_{\operatorname{Rep}(G)}
(F(M),W)$ where $F$ denotes the forgetful
functor
$$
F:\operatorname{Rep}(G)\longrightarrow
\operatorname{Vec}.
$$

\noindent
To show that $\operatorname{Coind}$ is right
adjoint to $F$ we have to show that the
association set up above is an isomorphism,
that is, that any linear map
$\overline{f}:M\to W$ yields an
$m\longmapsto\varphi_m$ as above. The
construction is going to be
$$
\varphi_m(g)=\overline{f}(gm).
$$
\end{proof}

\begin{proposition}
\label{proposition-adjoint-exact-4-sept}
An additive functor which is a right adjoint
is
left exact.
\end{proposition}

\begin{proof}
Exercise.
\end{proof}

\begin{proposition}
\label{proposition-monos-4-sept}
A left exact functor sends monomorphisms to
monomorphisms.
\end{proposition}

\begin{proof}
Exercise.
\end{proof}

\begin{proposition}
\label{proposition-adjoint-injective-4-sept}
The right adjoint to an exact functor sends
injective objects to injective objects.
\end{proposition}

\begin{proof}
Exercise.
\end{proof}

\begin{lemma}[Corollary]
\label{lemma-adjoint-injectives-4-sept}
If $A\to B$ is an exact functor which has a
right adjoint, and if $B$ has enough
injectives,
then $A$ has enough injectives.
\end{lemma}

\noindent
Proof of Lemma 3.4.4. The following paragraph
essentially replicates the proof of the
corollary above. Let
$$
\operatorname{Coind}:\operatorname{Vec}
\longrightarrow\operatorname{Rep}(G)
$$
denote the right adjoint to the forgetful
functor
$$
\operatorname{Rep}(G)\longrightarrow
\operatorname{Vec}.
$$
Here all vector spaces are over $F$, which we
omit from the notation. Let
$V\in\operatorname{Rep}(G)$ and consider its
image in $\operatorname{Vec}$. We want to map
it to some injective
$I'\in\operatorname{Vec}$.
Now, in $\operatorname{Vec}$ all objects are
injective, so $I'=V$ will do. Now define
$$
I=\operatorname{Coind}(I'),
$$
which by the adjunction receives a morphism
$V\to I$. The morphism is a monomorphism by
Proposition \ref{proposition-monos-4-sept}.

\noindent
The next question is whether the injective
$I$
in $\operatorname{Rep}(G)$ is injective in
$G$-mod. This is the content of the following
proposition.

\begin{proposition}
\label{proposition-coind-injective-4-sept}
Assume that $\operatorname{char}(F)=0$. For
every $F$-vector space $W$, the module
$$
\operatorname{Coind}(W)
=
\operatorname{Hom}_F(FG,W)
$$
is injective in $G$-mod.
\end{proposition}

\begin{proof}
We show that
$$
\operatorname{Hom}_{\mathbb{Z}G}
(-,\operatorname{Coind}(W))
$$
is exact as a functor from $G$-mod to
$\operatorname{Ab}$. By Exercise 3.2.5 this
suffices.

\noindent
For any $M\in G$-mod we consider
$$
\beta_M:
\operatorname{Hom}_{\mathbb{Z}G}
(M,\operatorname{Hom}_F(FG,W))
\longrightarrow
\operatorname{Hom}_{\mathbb{Z}}(M,W)
$$
defined by
$$
\varphi\longmapsto
\bigl[m\longmapsto\varphi(m)(1_{FG})\bigr].
$$
For its inverse we start with
$$
\operatorname{Hom}_{\mathbb{Z}G}
(M,\operatorname{Hom}_{\mathbb{Z}}
(\mathbb{Z}G,W))
\longleftarrow
\operatorname{Hom}_{\mathbb{Z}}(M,W)
$$
and the assignment
$$
\bigl[m\longmapsto
\bigl[x\longmapsto
\psi(x\mathbin{\cdot}m)\bigr]\bigr]
\longleftarrow\psi.
$$
Since $W$ is a vector space, any
$\mathbb{Z}$-linear map $\mathbb{Z}G\to W$
extends in a unique way to an $F$-linear map
$FG\to W$. So this inverse lands in the right
place, and a simple calculation confirms it
really is an inverse. Thus
$$
\operatorname{Hom}_{\mathbb{Z}G}
(-,\operatorname{Coind}(W))
\simeq
\operatorname{Hom}_{\mathbb{Z}}(-,W).
$$
In more sophisticated language, $\beta$
establishes a natural isomorphism between
these
two functors.

\noindent
But $\operatorname{Hom}_{\mathbb{Z}}(-,W)$ is
exact because $W$, being a vector space over
a
field of characteristic zero, is a divisible
abelian group, hence injective in
$\operatorname{Ab}$. This proves that
$\operatorname{Coind}(W)$ is injective in
$G$-mod.
\end{proof}

\begin{exercise}
\label{exercise-vec-unnecessary-4-sept}
The introduction of $\operatorname{Vec}$ in
addition to $\operatorname{Rep}(G)$ and
$G$-mod might seem unnecessary. Indeed let
$$
F:\operatorname{Rep}(G)\longrightarrow
G\operatorname{-mod}
$$
be the forgetful functor, and let
$$
G=\operatorname{Hom}_{\mathbb{Z}G}(FG,-).
$$
Is $G$ right adjoint to $F$? If so, then we
can
produce injective resolutions of
$V\in\operatorname{Rep}(G)$ by the corollary
above. Now the question is whether such
objects
are injective in $G$-mod.
\end{exercise}

\begin{exercise}
\label{exercise-adjoints-injective-4-sept}
Let $F:\mathcal{A}\to\mathcal{B}$ be an exact
functor, and let
$G:\mathcal{B}\to\mathcal{A}$
be a right adjoint to $F$. Is it true that
for
$I\in\mathcal{B}$ injective, the object
$F(G(I))\in\mathcal{B}$ is injective? The
answer is no, and a counterexample is given
by
the categories of modules over the rings
$\mathbb{C}[x]/(x^2)$ and $\mathbb{C}$.
\end{exercise}

\noindent
Find somewhere to discuss $i^*$ and its two
adjoints $i_*$ and $i_!$.

\begin{example}
\label{example-cohomology-c-c-star-4-sept}
Consider the short exact sequence of
$G$-modules
$$
0\longrightarrow\mathbb{Z}
\longrightarrow\mathbb{C}
\longrightarrow\mathbb{C}^{\times}
\longrightarrow0,
$$
where the first map is inclusion and the
second
is the exponential
$$
a\longmapsto e^{2\pi ia}.
$$
Applying Proposition 3.2.1, we obtain a long
exact sequence
$$
\cdots\longrightarrow
H^{n-1}(G;\mathbb{C}^{\times})
\longrightarrow H^n(G;\mathbb{Z})
\longrightarrow H^n(G;\mathbb{C})
\longrightarrow H^n(G;\mathbb{C}^{\times})
\longrightarrow H^{n+1}(G;\mathbb{Z})
\longrightarrow\cdots.
$$
By Lemma 3.4.4, $H^n(G;\mathbb{C})=0$ for
$n\geq1$. Therefore
$$
H^{n+1}(G;\mathbb{Z})
\simeq
H^n(G;\mathbb{C}^{\times})
$$
for all $n\geq1$.
\end{example}

\begin{exercise}
\label{exercise-cyclic-c-c-star-4-sept}
Confirm these isomorphisms in the case
$G=C_m$ using the general computation above,
applied to $A=\mathbb{Z}$ and
$A=\mathbb{C}^{\times}$.
\end{exercise}

\section{Example: Hilbert's theorem 90
4-sept}
\label{section-hilberts-theorem-90-4-sept}

\noindent
Recall that a finite field extension $L/K$ is
called Galois if $L^G=K$, where $G$ is the
Galois group
$$
\operatorname{Gal}(L/K)
=
\{f\in\operatorname{Aut}(L)mid f|_K
=\operatorname{Id}_K\}.
$$
The degree $[L:K]$ of the extension is the
dimension of $L$ as a vector space over $K$,
and it is a theorem that
$\#\operatorname{Gal}(L/K)$ equals the
degree.
Since $L$ is a finite dimensional vector
space
over $K$ we can consider, for each
$\alpha\in L$, the $K$-linear endomorphism
$L\to L$ of multiplication by $\alpha$ as an
element $m_\alpha\in\operatorname{GL}_K(L)$.
Then we define the norm of $\alpha$ to be
$$
N(\alpha)=\det(m_\alpha).
$$
It is a theorem that if $L/K$ is Galois then
$$
N(\alpha)=\prod_{g\in G}g(\alpha).
$$
We discuss the following result of Noether,
which is often called Hilbert 90.

\begin{theorem}[Noether]
\label{theorem-noether-hilbert-90-4-sept}
Let $L/K$ be a Galois extension, and let
$G=\operatorname{Gal}(L/K)$. Then $G$ acts on
$L$, and on $L^{\times}$, so that
$L^{\times}$ is a $G$-module. One has
$$
H^1(G,L^{\times})=0.
$$
\end{theorem}

\noindent
Hilbert only stated the theorem in the case
$G\simeq\mathbb{Z}/n$ is cyclic. The
statement is given as follows.

\begin{theorem}
\label{theorem-hilbert-90-cyclic-4-sept}
Let $L/K$ be a Galois extension, and suppose
that
$$
G=\left\langle{}\sigma\right\rangle{}
\simeq\mathbb{Z}/n.
$$
Then for every $\alpha\in L$ of norm
$N(\alpha)=1$, there exists $\beta\in L$ such
that
$$
\alpha=\sigma(\beta)/\beta.
$$
\end{theorem}

\noindent
We will give an elementary proof of Hilbert's
version first.

\begin{proof}
Let $\alpha\in L^{\times}$ with
$N(\alpha)=1$. Let $y\in L^{\times}$ and
consider
$$
d_0=\alpha y,
\qquad
d_{j+1}=\alpha\sigma(d_j),
\qquad
j=0,1,2,\ldots.
$$
Equivalently,
$$
d_j=
\bigl(\alpha\sigma(\alpha)\cdots
\sigma^j(\alpha)\bigr)\sigma^j(y).
$$
We note that
$$
d_{n-1}=N(\alpha)\sigma^{n-1}(y)
=\sigma^{n-1}(y).
$$
Let
$$
\beta=d_0+d_1+\cdots+d_{n-1}.
$$
We claim that there exists a choice of $y$
such
that $\beta\ne0$. For if there were not, we
would have
$$
\sum_{i=0}^{n-1}m_i\sigma^i(y)=0
$$
for all $y\in L^{\times}$. Thus there would
be
a linear dependence among the $\sigma^i(y)$
for
every $y\in L$. But by a certain theorem in
Galois theory we can choose a generator $y$,
that is, $L=K(y)$. Then the $n=[L:K]$ images
$$
\{\sigma^i(y)\mid\sigma^i\in G\}
$$
must form a basis, a contradiction.

\noindent
All right, so $\beta\ne0$. But then
$$
\begin{aligned}
\frac{\sigma(\beta)}{\beta}
&=
\frac{\sigma(d_0)+\cdots+\sigma(d_{n-1})}
{d_0+\cdots+d_{n-1}} \\
&=
\frac{d_1/\alpha+d_2/\alpha+\cdots+
d_{n-1}/\alpha+d_0/\alpha}
{d_0+d_1+\cdots+d_{n-1}} \\
&=\alpha^{-1}.
\end{aligned}
$$
Therefore
$$
\frac{\sigma(\beta^{-1})}{\beta^{-1}}
=
\frac{\beta}{\sigma(\beta)}
=\alpha.
$$
\end{proof}

\begin{exercise}
\label{exercise-hilbert-90-cyclic-4-sept}
Compare the statement for cyclic extensions
with the convenient complex for cohomology
of a
cyclic group, to understand why Hilbert's
theorem is a special case of the
cohomological
version of Noether's theorem.
\end{exercise}

\noindent
To prove Noether's theorem, we first need a
lemma. This proof, which I found on Noam
Elkies's website [2], is essentially from
Jacobson's book, though there it appears as
[?], more or less buried in the proof of
the lemma on page 236.

\begin{lemma}[Artin's lemma]
\label{lemma-artin-4-sept}
Let $L/K$ be a Galois extension of degree $n$
with Galois group
$G=\operatorname{Gal}(L/K)$. Denote by
$$
\rho:G\longrightarrow\operatorname{GL}_K(L)
$$
the canonical inclusion. If
$$
\sum_{g\in G}c_g\rho(g)=0
$$
for constants $c_g\in L$, then $c_g=0$ for
all $g$.
\end{lemma}

\noindent
Notice that $\operatorname{GL}_K(L)$ is a
$K$-vector space, so the lemma says the
$\rho(g)$ are linearly independent in this
vector space, but it says more: the
coefficients can be in $L$.

\begin{proof}
Let $\{x_1,\ldots,x_n\}$ be a $K$-linear
basis of $L$, and let $g_1,\ldots,g_n$ be the
elements of $G$ in some order. Consider the
matrix
$$
M\in\operatorname{Mat}_{n\times n}(L)
$$
defined by
$$
M_{ij}=g_i(x_j).
$$
Then the relation in the statement says
$cM=0$, where $c$ is the row vector
$$
(c_1,\ldots,c_n).
$$
This means $M$ must be a degenerate matrix,
so
there exist nonzero column vectors
$$
v=(v_1,\ldots,v_n)^T
$$
such that $Mv=0$. Clearly the set of all such
column vectors, denoted $V\subset L^n$, is an
$L$-vector space. For each $g\in G$ and
$v\in V$, the relation
$$
\sum_i v_i g(x_i)=0
$$
holds.

\noindent
We observe that $V$ is invariant under the
entry-wise action of $G$. Indeed if $v\in V$,
then
$$
\sum_i v_i h(x_i)=0
$$
for all $h\in G$, and hence
$$
g\left(\sum_i v_i g^{-1}h(x_i)\right)
=
\sum_i g(v_i)h(x_i)=0
$$
for all $g\in G$ too.

\noindent
We claim there exists such a $v\in V$ with
all
entries in $K$. Let $v\in V$ be nonzero with
a
a minimal number of nonzero entries. Without
loss
of generality $v_1=1$. Now if we apply some
$g$
to all entries we get another element of $V$,
and subtracting gives a vector
$$
w=v-g\ast v
$$
with first entry zero. By minimality of $v$
we
have $w=0$. But this means $g(v_i)=v_i$ for
each $g\in G$. Hence each $v_i\in L^G=K$.

\noindent
But with such a vector $v$ in hand, we set
$g=e$ to conclude that
$$
\sum_i v_ix_i=0,
$$
meaning the basis $\{x_1,\ldots,x_n\}$ would
not be a basis. This contradiction proves the
lemma.
\end{proof}

\noindent
The following proof of Noether's theorem can
be
found in [16, page 150].

\begin{proof}
Let $\alpha$ be a nonzero $1$-cocycle. Thus
$\alpha:G\to L^{\times}$ satisfies
$$
\alpha(gh)=\alpha(g)g(\alpha(h)),
$$
while a $1$-coboundary is such an $\alpha$
which is of the form
$$
\alpha(g)=\beta/g(\beta).
$$
We wish to show that every $1$-cocycle is a
$1$-coboundary.

\noindent
Consider
$$
\sum_{g\in G}\alpha(g)\rho(g).
$$
Since it is nonzero by Lemma
\ref{lemma-artin-4-sept}, there exists
$\theta\in L^{\times}$ such that
$$
\sum_{g\in G}\alpha(g)g(\theta)\ne0.
$$
Take $\beta$ to be this element. Then we
claim
that
$$
\alpha(g)=\beta/g(\beta),
$$
so we are done.
\end{proof}

\section{Example: Weyl's complete
reducibility theorem 4-sept}
\label{section-weyl-4-sept}

\noindent
We discuss cohomology of Lie algebras.
Although
some familiarity with Lie algebras would be
helpful, I think this discussion should be
comprehensible even to those without previous
exposure to the subject.

\begin{definition}
\label{definition-lie-algebra-4-sept}
A {\it Lie algebra} over a field $F$ is an
$F$-vector space $\mathfrak{g}$ equipped
with a
bilinear operation
$$
[\mathbin{\cdot},\mathbin{\cdot}]:
\mathfrak{g}\times\mathfrak{g}
\longrightarrow\mathfrak{g}
$$
satisfying
\begin{enumerate}
\item $[b,a]=-[a,b]$,
\item
$$
[a,[x,y]]=[[a,x],y]+[x,[a,y]].
$$
\end{enumerate}
\end{definition}

\noindent
A simple example of a Lie algebra, denoted
$\mathfrak{gl}(V)$, is
$\operatorname{End}(V)$ with
$$
[A,B]=AB-BA.
$$
Another is the set $\operatorname{Der}(R)$ of
derivations of a commutative $F$-algebra $R$.
A derivation is an $F$-linear map $D:R\to R$
satisfying
$$
D(ab)=D(a)b+aD(b).
$$
The Lie bracket is again
$$
[D_1,D_2]=D_1D_2-D_2D_1.
$$


codex resume 01a06df8-835b-7e23-b0a4-d2ecaba3a558

\bibliography{my}
\bibliographystyle{amsalpha}

\end{document}
